%% Generated by tools/build_guide.py from UNIT-REGISTRY.md.
%% Do not edit by hand: edit the source Markdown or the registry.
\chapter{编程指南}
\guidechaptermark{编程指南}
\label{ch:programming-guide}

\guidechapterintro{本章按功能重新组织同一门语言，说明 Croqtile 程序的结构、类型系统、符号形状、并行性、分块选择、数据搬运、收缩与协调，以及背后的设计依据。完整核函数就放在它所演示的特性旁边，因此每一节也可独立阅读。}

\section{程序结构与编译模型}
\label{sec:3-1}

一个 \co@.co@ 文件将应用程序与其中的并行工作放在同一份源码中。普通主机 C++ 准备输入并调用 \co@__co__@ 函数；后者描述并行工作，也可以调用设备辅助函数完成专门计算。理解每部分由哪个编译器处理，有助于代码集成和错误定位。


\subsection*{构造一个 Croqtile 程序}

一个典型的 \emph{Croqtile} 程序由几部分组成，具体取决于目标平台。

\begin{itemize}[leftmargin=1.6em]
\item 设备程序
\item 主机程序
\item tileflow 程序
\end{itemize}


\subsection*{转译与编译}

Croqtile 的编译过程通常包含三个主要步骤：\textbf{预处理}、\textbf{转译}和\textbf{目标编译}。

图~\ref{fig:f14-compile-workflow} 展示了这三个阶段。

\begin{figure}[htbp]\centering
\includegraphics[width=\linewidth,height=0.56\textheight,keepaspectratio]{f14-compile-workflow.drawio.png}
\caption{Croqtile 的完整编译流程}
\label{fig:f14-compile-workflow}
\end{figure}

如图所示，\emph{预处理}之后，Croqtile 立刻把 \emph{tileflow 程序} \emph{转译}（也就是\emph{源到源}编译）为目标代码，而用户提供的\emph{主机程序}和\emph{设备程序}保留原样。\emph{tileflow 程序}被\emph{转译}成主机源代码和设备源代码，这些统称 \emph{Croqtile 生成代码}。随后，编译器把用户代码和 \emph{Croqtile 生成}代码合并，执行\emph{目标编译}。这个过程可以产出多种结果，比如转译后的源代码、工作流脚本、目标模块、目标汇编和可执行文件。

Croqtile 驱动完整的构建过程，其中的关键步骤是把 \emph{tileflow 程序}转译成 \emph{Croqtile 生成代码}。

Croqtile 编译有一个显著特点：同时支持两种模型：

\begin{itemize}[leftmargin=1.6em]
\item \textbf{单源编译模型}：目标编译器允许在同一源文件中编译设备程序和主机程序。
\item \textbf{分离源编译模型}：主机代码和设备代码需要分别编译。
\end{itemize}

上面的代码天然支持\emph{单源编译模型}。但要支持\emph{分离源编译模型}，Croqtile 要求用 \co@__cok__@ 块把设备程序包起来，如下所示：

\begin{lstlisting}[language=croqtile,caption={转译与编译}]
__cok__ {
  void device_function(...) { ... }
} // end of __cok__
__co__ void croqtile_function(...) { ... }
void foo() { ... }
\end{lstlisting}

\co@__cok__@ 标记需要在调用目标编译器前从主机代码中分离的设备代码。当目标平台分别编译主机与设备代码时，需要使用这一包装。源码仍可维护在一个 \co@.co@ 文件中，由包装告诉构建过程在哪里拆分。


\subsection*{设备程序：并行计算}

设备辅助函数接收调用方选取并暂存的数据，用目标语言完成计算。下面的加法辅助函数展示这条边界：Croqtile 安排分块与搬运，辅助函数则使用传入的指针和元素数量。

下面再次列出前面的设备辅助函数：

\begin{lstlisting}[language=croqtile,caption={从 tileflow 调用的设备辅助函数}]
__device__ void kernel(int * a, int * b, int * c, int n) {
  for (int i = 0; i < n; ++i) c[i] = a[i] + b[i];
}
\end{lstlisting}

对于只支持\emph{分离源编程模型}的目标平台，代码可能需要包在 \co@__cok__@ 块里。等价的写法如下：

\begin{lstlisting}[language=croqtile,caption={分离源所需的 cok 包装}]
__cok__ {
  extern "C" void kernel(int * a, int * b, int * c, int n) {
    for (int i = 0; i < n; ++i) c[i] = a[i] + b[i];
  }
} // end of __cok__
\end{lstlisting}

这里 \co@extern "C"@ 标注取代了目标平台使用的 \co@__device__@ 关键字，因为设备函数要求 C 链接。

设备程序遵循单程序多数据范式。主机与设备之间、以及设备内部各存储层次之间的数据搬运，由 Tileflow 程序统一编排；设备函数则负责具体计算。


\subsection*{异构性：从转译到程序启动}

在某些目标平台上，\co@parallel-by@ 块不仅指定并行线程的数量，还划出了主机代码与设备代码的边界，这正是 Croqtile 管理异构性的方式。具体来说：

\begin{itemize}[leftmargin=1.6em]
\item \co@parallel-by@ 块\textbf{之外}的代码被转译为\textbf{主机代码}。
\item 块\textbf{之内}的代码被转译为\textbf{设备代码}，如下图所示：
\end{itemize}

\begin{lstlisting}[language=croqtile,caption={异构性：从转译到程序启动}]
__co__ void foo(...) {
  // Generate host code
  // ...
  parallel p by 6 {
    // Generate device code
    // ...
  }
  // Generate host code
  // ...
}
\end{lstlisting}

熟悉相关目标平台的读者会看出，这里正对应着 \co@parallel-by@ 语句处发生的程序启动。Croqtile 用统一的单程序多数据并行模型把这些细节抽象掉，让开发者专心写算法，不必操心底层异构性。


这份参考覆盖 Croqtile 编程的方方面面：

\begin{itemize}[leftmargin=1.6em]
\item \textbf{程序结构}：Croqtile 程序如何组织——\co@__co__@ 函数、主机侧 API 与 C++ 互操作。
\item \textbf{带形状数据}：Croqtile 面向多维张量的类型系统——形状、\co@MDSpan@、\co@spanned data@、符号维度与动态形状。
\item \textbf{循环与并行性}：单程序多数据并行、循环控制、分块与迭代构造。
\item \textbf{数据搬运}：DMA 语句、用 \co@chunkat@ 分块、重排与高级数据搬运模式。
\item \textbf{C++ 嵌入}：输入／输出约定、\co@__cpp__@ 内联块、宏与预处理器集成。
\item \textbf{优化模式}：tileflow 优化、异步 DMA 模式、多缓冲与性能调优技术。
\end{itemize}


\subsection*{小结}

到这里可以看清：Croqtile 程序由三部分组成，其中 \emph{tileflow 程序}是核心，它在编译时被转译成目标源代码。调用链通常从主机到 tileflow 程序，再从 tileflow 程序到设备代码。

Croqtile 不仅把 DMA 操作抽象成更高层的语义，还把迭代和分块结合起来，因此代码既简洁又有表现力。



\section{语法设计与类型系统}
\label{sec:3-2}

\subsection*{类型系统补记：\texttt{mdspan} 与 \texttt{ituple}}

选择一个分块需要知道张量形状、各轴划分为多少块，以及当前块的坐标。Croqtile 用形状、元组和有界索引显式表示这些信息，让分块表达式保留原本分散在偏移计算中的关系。

带形状声明 \co@s32 [128, 256] lhs@ 同时确定元素类型和秩。这里的尺寸是常量，编译器能在执行前求值，拒绝用三个坐标访问二维张量，并检查分块形状的兼容性。声明使用符号维度时，即使部分尺寸来自调用方，秩信息仍然保留。

\co@span(i)@ 语法取出一个维度：\co@lhs.span(0)@ 是 128，\co@rhs.span(1)@ 是 256。简写 \co@lhs.span@ 则复制整个形状。

\co@ituple@ 将整数值组成元组。在 \co@parallel {px, py} by [8, 16]@ 中，坐标还带有范围：当前值标识某个实例，上界描述整个网格。编译器可以化简上界及派生表达式，但并不是给所有实例确定同一个当前索引。依赖实际执行的索引计算仍在运行时进行。

第 2.4 和 2.5 节介绍具体声明。这里需要区分程序本身已知的信息与某次调用提供的信息。编译时检查使用前者，入口检查和运行时断言则处理依赖后者的约束。





它的目标是给日常编写高性能核函数的工程师减负。借助 Croqtile，编程过程更容易上手，工程师也能更轻松、更高效地优化这些关键环节。


\subsection*{把 Croqtile 嵌入 C++}

Croqtile 是嵌入在 C++ 里的领域专用语言代码。Croqtile 编译器对 Croqtile 函数做源到源翻译，生成 C++ 代码，并隐式包含 \co@choreo.h@，让翻译后的代码能与其他 C++ 代码正常协作。下面的代码片段是一个例子。

\begin{lstlisting}[language=croqtile,caption={在 C++ 中嵌入 Croqtile}]
// some C++ code
__co__ void croqtile_function() {
  // croqtile code
}
// another C++ code
void foo() {
  croqtile_function();
}
\end{lstlisting}

\co@__co__@ 前缀标识需要由 Croqtile 翻译的函数体。普通 C++ 通过声明的函数名调用它。面向主机的接口使用带形状的缓冲区包装对象，使调用同时携带尺寸和地址。

此外，Croqtile 引入了一些专用类型，用来保证 Croqtile 函数的实参与调用方之间的一致性。相关细节将在数据类型的讨论中说明。


\subsection*{变量与类型}

Croqtile 引入四类基础类型：标量类型、带形状类型、整数元组类型（\co@ituple@ 类型）和有界类型。它们分别服务于不同的用途。

\begin{itemize}[leftmargin=1.6em]
\item \textbf{标量类型}：为满足程序控制的需要而设计，由\emph{整数类型}和\emph{布尔类型}组成。
\item \textbf{带形状类型}：表示参与计算的数据类型（通常是张量）。除了引用原始数据，\emph{带形状类型}还把数据与一个用多维范围表示的形状关联起来，这种「多维范围」就叫 \co@mdspan@。它对分块等操作很有用。
\item \textbf{整数元组类型}：表示一组整数值，常见用法是给出多维数据的界。
\item \textbf{有界（整数／ITuple）类型}：这是一类特殊类型，用于简化对数据子区域的引用，并与循环构造配合，迭代地处理数据。
\end{itemize}

四类之中，\emph{标量}和 \emph{ituple} 比较容易理解，因为它们对应现有通用编程语言里就有的要素；\emph{带形状类型}和\emph{有界类型}则是这个领域特有的。


\subsection*{带形状类型}

带形状类型是一种\textbf{复合类型}，由\textbf{基础类型}和**多维跨度（\co@mdspan@）类型**组成。\emph{基础类型}和 \emph{mdspan} 都不是\emph{完整类型}，也就是说，任一部分都不能单独声明带存储的数据。在 Croqtile 中，称它们为\emph{部分类型}。

将类型分为两部分，是为了能单独操作 \emph{mdspan}。机器学习数据通常按多维方式组织，程序员经常要把它切成更小的分块，让分块跨存储层次移动。因此，形状计算可以独立于元素类型进行。

部分类型必须组合成完整类型，才能声明带存储的多维数据。


\subsection*{部分类型：\textbf{mdspan}}

与多数类型系统不同，\emph{mdspan} 作为一种类型手段并不完整。称它为\emph{部分}，是因为 Croqtile 程序不能把数据定义成 \emph{mdspan} 类型，但 \emph{mdspan} 本身可以单独定义。这样设计基于一个观察：循环分块只关心多维数据的形状。因此，Croqtile 允许程序员独立于基础类型操作 \emph{mdspan}，从而简化形状计算。

与此同时，Croqtile 编译器可以对 \emph{mdspan} 做类型检查，尽可能早地（在适用时早于执行）暴露代码错误。

在 Croqtile 中，mdspan 用 \co@[@ 和 \co@]@ 括起来，其中的维度上界用 \co@,@ 分隔：

\begin{lstlisting}[language=croqtile,caption={一个 mdspan 声明}]
mdspan sp : [7, 8];           // defined a mdspan of 2-dimensions
mdspan<4> mds : [a, 3, 4, 1]; // 'a' is an existing integer
d : [c, 4, 28];               // 'd' is not explicit annotated. Type is deduced.
\end{lstlisting}

如上例所示，可以用 \emph{mdspan} 部分类型关键字显式定义 mdspan，也可以选用 \co@<@ \co@>@ 括起的秩。还可以不加类型标注定义 mdspan，比如上面 \co@d@ 的定义。（注意：\co@mds@ 和 \co@d@ 属于\textbf{依赖类型}，因为它们的类型取决于 \co@a@ 和 \co@b@ 的求值。程序员可以给 \co@a@、\co@b@ 提供运行时值，但这样一来就要付出一些运行时检查的代价。）

对 C++ 程序员来说，可以把 mdspan 看成多维数组的特征描述，它描述的是多层范围。例如上例中的 \co@sp@ 定义了两层范围，分别是 0 到 6 和 0 到 7。

定义好 mdspan 之后，可以用 \co@()@ 运算符取出整数维度值。

\begin{lstlisting}[language=croqtile,caption={一个形状字面量}]
sp : [7, 8];
int b = sp(0) + sp(1);   // 'b' equals to 15 (7 + 8)
\end{lstlisting}

上例中，表达式 \co@sp(0)@ 求出 mdspan \co@sp@ 第一维的取值。维度值本质上是整数，可以直接参与整数算术。

于是就可以像下面这样，从已有的 mdspan 派生出新的 mdspan：

\begin{lstlisting}[language=croqtile,caption={改变一个维度后的形状}]
sp : [6, 8];
spn : [1, sp(0)/2, sp(1)/4]; // define a new mdspan from the existing one.
\end{lstlisting}

有了这个机制，在 Croqtile 代码里对多维跨度做分块更加方便。分块是高性能核函数中的常见操作，因此 Croqtile 提供了相应的语法糖：

\begin{lstlisting}[language=croqtile,caption={同样的形状，另一种写法}]
sp : [6, 8];
spn : sp [1, (0)/2, (1)/4];  // spn is defined as [1, 3, 2]
\end{lstlisting}

这段代码和前面那段作用相同，但语法更简洁。这体现了 Croqtile 的一条设计理念：用简洁的代码表达常见操作。

命名 \co@mdspan@ 必须在声明时初始化，之后不能修改。需要另一种形状时，应从已有形状派生新名字，使编译器能够保留并分析形状表达式之间的关系。

除了这种\emph{逐维度}定义 mdspan 的方式，Croqtile 还支持其他定义方法；这些方法在 \emph{i-tuple} 的讨论中介绍。


\subsection*{有界类型}

有界整数同时携带当前值和排他上界。在 \co@parallel p by 6@ 中，可能的取值是 0 到 5，上界为 6。因此，分块选择既能确定 \co@p@ 标识哪一块，也能确定总共有多少块。\co@parallel-by@ 和 \co@with-in@ 建立这种关联。

有界元组将同样的做法扩展到多个维度。每个分量有自己的上界，整个范围由 \co@mdspan@ 描述。例如，范围为 \co@[4, 8]@ 的元组能标识 32 个位置之一。第 2.7 节说明如何声明和遍历这些索引。





\subsection*{小结}

Croqtile 以嵌入 C++ 的领域专用语言表达单程序多数据程序，保留 C++ 风格的写法。它主要简化通过 DMA 跨存储层次搬运数据的底层细节，因此也被称为数据流编程语言。其目标是支持日常核函数开发，让开发者更专注于计算组织，减少重复的实现工作。

实际使用中的反馈，也有助于进一步改进这些构造。


\subsection*{范围的组与整数的组}

在 Croqtile 中，\emph{mdspan} 可以表示各维度的形状。\emph{mdspan} 就是\textbf{多维跨度}，其中的元素是\emph{维度跨度}，隐含着一个取值范围。

例如，\emph{mdspan} \co@[7, 14]@ 表示两个范围组成的一组：\co@[0, 7)@ 和 \co@[0, 14)@。所以 \emph{mdspan} 并不表示一组整数值。不过在某些场景里，把整数值组成一组确实有用：

\begin{itemize}[leftmargin=1.6em]
\item 计算两个 mdspan 之差时需要多个整数；在高层的语义上，这个差可能反映调整形状维度时的填充差异。
\item 对形状做分块时，各个分块因子是整数，而不是范围。
\end{itemize}

这些情形促使 Croqtile 提供一种把整数值组成一组的类型。


\subsection*{前瞻：有界 I-Tuple}

有界元组为普通整数元组附加范围。例如，位置 \co@{1, 2}@ 与范围 \co@[4, 8]@ 结合后，既标识当前位置，也描述容纳它的空间。循环用范围确定迭代，\co@chunkat@ 用范围推导划分。


\subsection*{\emph{Spanned}：数据与缓冲区}

tileflow 程序的主要任务，就是搬动大规模数据集来处理它们。在某些术语体系里，Croqtile 函数的输入和输出分别称为「输入数据」和「输出数据」，其余存储则称为「缓冲区」。从函数的视角看，这就在概念上分出了「外部」内存和「内部」内存。不过「数据」和「缓冲区」指的都是存储位置；在 Croqtile 中，两者都定型为\textbf{带形状类型}，表明它们是与形状关联的数据或缓冲区。

这两种情况都称为带形状数据。\co@mdspan@ 描述形状，带形状对象则将形状、元素类型和存储访问组合起来。单独命名或推导一个形状，并不会复制数据。


\subsection*{缓冲区生存期管理}

在 \emph{Croqtile 函数}内部分配的存储，生存期由 Croqtile 编译器管理，以便高效利用：只要生存期不重叠，编译器就尽量复用缓冲区。因此，开发者不必自行管理这些缓冲区。

自动分配不能替代同步。若异步搬运或计算仍在使用缓冲区，必须等它结束后才能覆盖或释放存储。流水线中的等待和完成操作明确了这些使用时段，使程序能够安全复用同一组槽位。


\subsection*{有界变量}

在 \co@parallel-by@ 块中，紧跟 \co@parallel@ 关键字定义的变量称为\textbf{有界变量}。具体来说：

\begin{itemize}[leftmargin=1.6em]
\item 若变量是\emph{整数}，称为\textbf{有界整数}。
\item 若变量是 \emph{ituple}，称为\textbf{有界 ituple}。
\end{itemize}

这个术语表明，变量不只有值，还有确定的界。例如在语句 \co@parallel p by 6@ 中，\co@p@ 是\emph{有界整数}，界为 \co@[0, 6)@，其中 \co@6@ 是不含在内的上界。

并行索引的当前值标识执行实例，循环索引的当前值随 \co@foreach@ 迭代变化。两者的上界分别来自 \co@parallel-by@ 和 \co@with-in@ 声明。


\subsection*{有界变量的取值}

有界变量同时携带范围和当前值。对于 \co@with-in@ 声明的索引，\co@foreach@ 按以下规则提供当前值：

\begin{itemize}[leftmargin=1.6em]
\item \textbf{当前值}：在 \co@foreach@ 语句中，有界变量充当迭代变量，当前值由循环迭代决定；在 \co@foreach@ 语句之外，其当前值恒为零。
\end{itemize}

\begin{itemize}[leftmargin=1.6em]
\item \textbf{上界值}：由 \co@with-in@ 或 \co@parallel-by@ 语句指定。
\end{itemize}

\co@foreach@ 结束后，\co@with-in@ 索引恢复为循环外的值，不会像 C++ 循环变量那样保留终止值。并行索引遵循另一条规则，其值标识当前并行实例。下面的例子讨论顺序执行的 \co@with-in@ 情形。

\begin{lstlisting}[language=croqtile,caption={有界变量的取值}]
with x in 6 {
  // x's current value is 0
  foreach x {
    // x's current value is either 0, 1, 2, ..., 5
  }
  // x's current value is 0, NOT 6
}
\end{lstlisting}


\subsection*{Croqtile 中的内存}

当不同并行实例需要以不同方式使用分块时，存储安排便成为算法的一部分。同一个 \co@[M, N]@ 形状可能用于全局输入、共享暂存分块或私有工作缓冲区。形状控制索引，存储说明符确定缓冲区所在的位置。


\subsection*{Croqtile 中的带形状类型}

\textbf{带形状类型}是一种复合类型，同时包含基本数据类型（如 \co@f32@、\co@s32@、\co@f16@）和描述数据形状与组织的\textbf{多维跨度（\texttt{mdspan}）}。有了它，Croqtile 能用高效的布局定义大型数组或矩阵，程序员也可以抛开数据类型单独操作形状（也就是各个维度）。并行计算和机器学习负载里常见的多维数据集，正需要这种灵活性。

带形状类型本身并不为数据分配存储。它描述数据在内存中如何组织，从而为运行时留出灵活性与优化空间。


\subsection*{Croqtile 中的存储说明符}

在这里使用的目标平台上，\co@global@、\co@shared@ 和 \co@local@ 分别表示设备全局存储、块内共享存储和线程私有存储。第 2.6 节给出声明示例。先根据哪些实例需要读写缓冲区选择位置，再按需安排不同位置之间的搬运。



\section{符号形状、形状求值与值编号(value numbering)}
\label{sec:3-3}

\subsection*{一等公民：形状}

仅凭一个原始指针，调用方无法知道元素组成的是向量、矩阵，还是从大张量中切出的分块。计算偏移或检查搬运时，代码就必须从别处取得这些信息。Croqtile 将形状与数据放在一起，使声明、分块选择和搬运语句能引用同一组尺寸。

Croqtile 则要求\textbf{任何声明或使用的数据都必须关联一个形状}，以此强化代码安全，也让带形状数据的编程变简单。这就是 Croqtile 把形状当作一等公民的原因。


\subsection*{\emph{mdspan} 的求值}

在很多场景里，取常量值的 mdspan 就足够了，因为高性能设备核函数往往要针对固定的输入形状做精细调优。这时 \co@mdspan@ 在编译时求值，不会带来额外的执行时间或存储开销。

当维度来自输入时，生成的主机代码在进入 Croqtile 函数时求值所需的形状表达式。常量部分仍可提前化简，运行时工作量取决于余下的表达式与检查。


\subsection*{整数与 I-Tuple 的求值}

操作数已知时，整数和元组表达式在编译时求值；依赖运行时索引或输入尺寸的表达式，则保留为生成程序中的计算。不可变性有助于编译器复用和化简表达式，但不会使每个值都成为常量。


核函数往往需要在保持程序结构的同时接受不同输入尺寸。匿名维度让大小暂时未知，符号维度则进一步为它命名。当两份输入的尺寸必须一致，或输出维度来自输入时，这个名字就很有用。


\subsection*{固定形状与动态形状}

固定形状的数据在高性能计算核函数里很常见，因为维度已知，可以做激进的优化。但有些场景必须处理带运行时维度的形状，也就是动态形状；这就要求核函数足够通用，能应对形状各异的输入。许多机器学习框架都强调这一点。


\subsection*{代码安全的提升}

使用\emph{符号维度}能提升代码安全，是因为 Croqtile 编译器可以施加更全面的检查。例如，先看 \co@Matmul@ 函数的\emph{匿名维度}版本：

\begin{lstlisting}[language=croqtile,caption={使用匿名维度的核函数}]
__co__ auto Matmul(s32 [?, ?] lhs, s32 [?, ?] rhs) {
  int tile_m = 32, tile_n = 8, tile_k = 16;
  s32 [lhs.span / {tile_m, tile_k}] tiled_lhs;  // 'mdspan' and 'ituple': rank must be same
  s32 [rhs.span / {tile_k, tile_n}] tiled_rhs;  // 'tiled_rhs' can not have a zero dimension
  // ...
}
\end{lstlisting}

Croqtile 同时做编译时检查和运行时检查来保证安全：

\begin{itemize}[leftmargin=1.6em]
\item \textbf{编译时检查}：Croqtile 校验秩是否一致。本例中 \co@lhs.span@ 的秩为 \co@2@，与 \co@tile_m@、\co@tile_k@ 的秩一致，所以没有问题。
\item \textbf{运行时检查}：Croqtile 生成的代码在运行时校验维度。例如 \co@tiled_lhs@ 的声明里，形状不能出现 \co@0@ 维，否则会得到非法的零大小缓冲区。Croqtile 的\emph{转译}过程会生成如下\emph{目标主机代码}：
\end{itemize}

\begin{lstlisting}[language=croqtile,caption={转译后的函数签名}]
void __croqtile_transpiled_Matmul(const croq::spanned_view<croq::s32, 2> & lhs,
                                  const croq::spanned_view<croq::s32, 2> & rhs) {
  croq::runtime_check(lhs.shape()[0] / 32 > 0,
                        "the 0 dimension of 'lhs' may result in spanned data "
                        "with a dimension value of 0.");
  // ...
}
\end{lstlisting}

条件不满足时，\co@croqtile_abort@ 函数会中止程序。这项前置检查帮助 Croqtile 及时发现问题。

对于\emph{符号维度}版本，Croqtile 编译器会在\emph{目标主机代码}里生成额外的检查：

\begin{lstlisting}[language=croqtile,caption={生成的形状检查}]
void __croqtile_transpiled_Matmul(const croq::spanned_view<croq::s32, 2> & lhs,
                                  const croq::spanned_view<croq::s32, 2> & rhs) {
  // additional check happens only when using symbolic dimensions
  croq::runtime_check(lhs.shape()[1] == rhs.shape()[0],
                        "dimension 'K' is not consistent.");
  croq::runtime_check(lhs.shape()[0] / 32 > 0,
                        "the 0 dimension of 'lhs' may result in spanned data "
                        "with a dimension value of 0.");
  // ...
}
\end{lstlisting}

除了\emph{匿名维度}版本已有的检查，\emph{符号维度}版本还会校验 \co@lhs@ 与 \co@rhs@ 之间维度是否一致，所以代码更安全。

因此建议使用符号维度：不只是为了易用性，也是为了代码更安全。


\subsection*{代价与限制}

为了支持动态形状，Croqtile 在编译时必须保守。这一限制让某些编译时检查做不了，只能挪到运行时。实现上已经尽量把这些检查放进主机代码来压低开销，但它们\textbf{仍然比固定形状的代码更保守}。程序员需要考虑这一限制，必要时添加用户级断言，加强代码安全。

\emph{匿名维度}和\emph{符号维度}都有一个重要限制：它们**只能用在参数列表里的 \emph{spanned data} 上**。内部缓冲区可以从输入派生形状，但不能用 \co@?@ 或新的符号名引入尚未指定的维度。例如，下面的代码会产生编译时错误：

\begin{lstlisting}[language=croqtile,caption={代价与限制}]
__co__ void foo(s32 [M] input) {
  s32 [M] derived;  // ok: the shape comes from the parameter
  s32 [?] fresh;    // error: an internal buffer cannot introduce an anonymous dimension
}
\end{lstlisting}


\subsection*{小结}

匿名维度表达未知大小，符号维度还使程序能够命名与之有关的关系。两个操作数共用 \co@K@，就记录了收缩维度必须一致的要求。Croqtile 在编译时检查能够确定的部分，并为依赖实际输入尺寸的关系生成检查。





\subsection*{\textbf{a. 可读性与清晰度}}

在形状定义里显式写出 \co@M@、\co@N@ 这样的符号维度，程序会\textbf{更易读}、也更好懂。读者不用看到 \co@?@ 或 \co@-1@ 再猜这个维度代表什么，而能立刻明白 \co@M@ 和 \co@N@ 是依赖运行时值的变量。这提升了\textbf{可维护性}，也减轻了阅读和修改代码的认知负担。

例如，两个问号表示两个未知尺寸，并没有声明两者相等。重复使用同一个符号名，则能直接在函数签名中记录相等关系：

\begin{lstlisting}[language=croqtile,caption={a. 可读性与清晰度}]
// Using symbolic dimensions
mdspan sp : [M, N];
mdspan spn : sp [1, M / 2, N / 4];
// Clear how the new shape is derived
// Using `?` or `-1`
mdspan sp : [? , ?];
\end{lstlisting}

维度相互独立时，匿名形式仍然有用。若维度之间的关系是操作要求的一部分，就应为它们命名。


\subsection*{\textbf{b. 类型安全与编译时检查}}

符号维度让编译器不必知道每个尺寸的数值，也能推导表达式。例如可以检查秩，并化简重复出现的 \co@M / 2@ 等关系。但这不会消除运行时验证：调用方提供的尺寸是否一致、是否满足所需约束，可能只有进入函数时才能确定。

下面的表达式将派生形状与命名的输入维度关联起来。编译器可以使用这项关联，读者也能直接看到它，不必再通过独立的尺寸参数维护关系：

\begin{lstlisting}[language=croqtile,caption={b. 类型安全与编译时检查}]
mdspan sp : [M, N];
// M and N are symbolic dimensions
mdspan spn : sp [1, M / 2, N / 4];
// Compiler ensures validity of shape expressions`
\end{lstlisting}

这保证了类型一致性，也减少了运行时出错的可能。


\subsection*{\textbf{c. 表达力与灵活性}}

符号维度描述数据形状的方式\textbf{更丰富}。\co@?@ 或 \co@-1@ 这类任意占位符只能表达基本的维度信息，符号维度却能用复杂的算术和关系来\textbf{约束}。比如，某个维度如何依赖另一个维度、如何缩放或分块，都可以直接写出来，自然得多。

匿名维度仍然可以通过输入形状读取。当同一个值出现在多个表达式中时，符号名更方便：修改某项推导后，它与其他维度之间的关系仍然清晰可见。

例如，用符号变量可以轻松表达维度之间的关系：

\begin{lstlisting}[language=croqtile,caption={c. 表达力与灵活性}]
mdspan sp : [M, N]; // Symbolic dimensions
mdspan spn : sp [1, M / 2, N / 4]; // Derived shape based on symbolic expressions
\end{lstlisting}

这种写法把依赖关系放在形状声明处。每次调用仍需满足由此产生的尺寸和整除要求。


\subsection*{\textbf{d. 编译时优化}}

符号表达式让编译器有机会复用维度计算，并根据已知关系化简表达式。常量部分可以在编译时求值，依赖输入大小的部分仍需运行时计算。符号名为优化保留了信息，但不会把未知输入尺寸变为常量，也不保证某种性能结果。


\subsection*{Croqtile 值编号简介}

在 Croqtile 中，值编号过程用来支撑形状推导：要么直接推导出编译时已知的形状，要么生成形状计算表达式。

许多值编号方法着眼于消除冗余计算，也就是那些始终产生相同值的表达式；Croqtile 的值编号则想建立一种渐进的方法，化简自己的领域专用构造 \co@mdspan@。在这种语境下，常量传播、代数化简等技术更为关键。

实现上，Croqtile 采用一种简单的\textbf{基于作用域}的值编号方法：它是局部值编号的扩展，配上一张词法作用域的值编号表[1]。与其他方法不同，Croqtile 不需要构造基本块、剪枝 SSA，也不需要编译器后端，因为它直接把值编号施加在抽象语法树上。但为了配合这种做法，必须给 Croqtile 的语法规则加上若干限制，尤其是防止值被改写。


\subsection*{快速示例}

值编号过程会给所有含有一组取值的实体分配值编号，mdspan 和 ituple 都在其中。例如：

\begin{lstlisting}[language=croqtile,caption={一个形状表达式}]
mdspan a : {1, b, c + 3}
\end{lstlisting}

值编号过程会产生下面这些取值：

\begin{lstlisting}[language=croqtile,caption={对应的值编号图}]
   VN #0:  const_1
   VN #1:  b
   VN #2:  c
   VN #3:  const_3
   VN #4:  +:#2:#3
   VN #5:  #0,#1,#4
\end{lstlisting}

图中「VN \#」后面的数字，是某个表达式的值编号。它可能表示编译时已知的常量整数，也可能表示符号值（如变量 \co@a@、\co@b@），还可能表示其他值编号的组合。这里值编号 \#0 和 \#3 引用常量值，\#1 和 \#2 分别引用整型变量 \co@b@ 和 \co@c@。\#4 表示一个加法表达式的值，也就是 \#2 与 \#3 之和；把所有值编号解引用之后，它就是表达式 \co@c + 3@。最后，值编号 \#5 表示由 mdspan \co@a@ 以符号方式定义的一组取值。有了这样的定义，就可以翻译索引操作了：例如 \co@a(0)@ 得到值编号 \#0，也就是常量 1。





\subsection*{从表达式到值编号}

在当前实现中，值编号以语法制导的方式进行：遍历抽象语法树的节点，遇到含表达式的节点就生成值编号。

生成值编号的方法很简单：把表达式作为输入，返回它对应的值编号。

更具体地说，它：

\begin{itemize}[leftmargin=1.6em]
\item 先由表达式生成一个签名。
\item 如果这个签名已经存在，就查询值编号表，返回与它关联的值编号。
\item 否则，在值编号表里把这个签名和一个新值编号关联起来，再返回这个新编号。
\end{itemize}

为简化实现，目前签名就是表达式的字符串；以后可以像多数值编号过程那样改用哈希来改进。

此外，为了处理符号的定义和引用，它遵循以下规则：

\begin{itemize}[leftmargin=1.6em]
\item 表达式用于初始化某个符号（变量或部分类型）时，该符号获得与表达式相同的值编号。更具体地说，它把\emph{带作用域}的符号名签名和值编号关联到表中。
\item 表达式引用某个符号时，就在值编号表里查找对应的签名并返回。
\end{itemize}

注意，生成签名时，符号必须使用\emph{带作用域的名字}，因为不同作用域里允许定义同名符号。


\subsection*{化简}

值编号过程每生成一个签名，都会尝试化简表达式。例如：

\begin{lstlisting}[language=croqtile,caption={一个常量表达式}]
   int a = 3 + 8;
\end{lstlisting}

对这个简单的整数表达式，值编号过程可能生成下面的签名和关联值编号：

\begin{lstlisting}[language=croqtile,caption={未化简的图}]
   VN #0:  const_3
   VN #1:  const_8
   VN #2:  +:#0:#1
   Alias a -> #2
\end{lstlisting}

如果不做优化，这里会生成两个常量值 \#0 和 \#1，以及一个表示两者相加的新值编号 \#2，它的签名是 \co@+:#0:#1@。但常量相加可以在编译时求值，所以可以化简。下面的输出展示了优化后的值编号过程：

\begin{lstlisting}[language=croqtile,caption={折叠后的图}]
   VN #0:  const_3
   VN #1:  const_8
   VN #2:  const_11
   Alias a -> #2
\end{lstlisting}

值编号 \#2 的签名现在表示一个常量，说明代数化简已经生效。


\subsection*{从值编号回到表达式}

经过值编号和化简，每个表达式都关联到一个化简后的值编号。在 Croqtile 中，这是 \co@mdspan@ 类型形状推导的重要一步。

\emph{mdspan} 类型是一种部分类型，定义基础类型元素的个数。编程时，它和某个基础类型一起，决定一个 \emph{spanned} 对象占多少内存。

\emph{mdspan} 的形状细节可以在运行时确定。但很多场景下，形状能在编译时完全求成固定数值。这时还可以在编译时触发静态内存分配，充分利用系统内存；否则就得生成计算形状的代码。

所以，除了常量值，还需要从值编号还原表达式。





\subsection*{词法作用域与限制}

在 Croqtile 中，值编号是\textbf{全局地}施加在抽象语法树上的。这是语法制导的方法，没有基本块，也就没有数据流分析，因此无法推导维度会变化的形状。

阅读生成代码时，可以看到直接结果：常量形状可能已折叠成常数，符号形状则成为生成程序中的计算。值编号尽量保留并化简这些关系，但不会证明所有数学恒等式。因此，手工看来等价的表达式，有时仍需要改用更明确的写法。


\subsection*{限制}

更具体地说，有两条限制：

\begin{itemize}[leftmargin=1.6em]
\item 任何 mdspan 都不能含有有界整数类型的元素，也不能由有界 ituple 生成。
\item 简单整数类型和 ituple 类型不能被重新定义。
\end{itemize}


\subsection*{参考文献}

[1] Briggs P, Cooper K D, Simpson L T. Value numbering[J]. Software: Practice and Experience, 1997, 27(6): 701-724.



\section{并行性、空间说明符与加掩码(masking)}
\label{sec:3-4}

\subsection*{并行性：Croqtile 如何表示并行工作}

前两章使用了较简单的例子：\co@parallel {i, j} by [4, 8]@ 创建 32 个实例，每个实例处理一个元素，却尚未指定这些实例的执行位置。对逐元素加法和理解数据搬运来说，这样处理没问题——编译器会选择默认配置。


\subsection*{如何看待并行任务}

在接触 Croqtile 语法之前，先把「并行」到底是什么意思说清楚。

假设有八个彼此独立的任务——八个要相加的分块、八行要处理的数据，等。你可以一次把八个任务都描述出来，称之为「并行」。但这句话完全没说明有多少任务真的在同时运行：在单核 CPU 上，它们一个接一个执行；在 4 核 CPU 上，也许同时执行四个。

\begin{figure}[htbp]\centering
\includegraphics[width=\linewidth,height=0.56\textheight,keepaspectratio]{f05-fig1_virtual_parallelism_light.png}
\caption{同样的八个任务，在每次执行一个、四个或八个任务的示意调度下，分别需要八个、两个或一个时间步。}
\end{figure}

并行性描述哪些任务可以独立执行，执行资源则决定同一时刻能推进多少任务。

如果分块数量和启动尺寸由同一组常量表示，修改分块就会迫使资源映射一起改变。将两种描述分开后，这些选择更容易调整。

Croqtile 用 \co@parallel@ 描述独立工作，用空间说明符描述映射。因此，可以先理解工作划分，再判断某个轴应映射到块、线程还是协作组。


\subsection*{矩阵乘法：组合起来}

有了 \co@parallel@、\co@foreach@、\co@dma.copy@、\co@chunkat@、\co@#@ 和空间说明符，分块矩阵乘法便可以组合实现。

思路是：用 \co@[128, 256]@ 的矩阵乘以 \co@[256, 256]@ 的矩阵，得到 \co@[128, 256]@ 的结果。


\subsection*{控制结构}

Croqtile 表达并行和循环的方式与 C++ 不同，但程序内部的分支处理仍然沿用 C++ 的 \co@if@-\co@else@ 块。


\subsection*{顺序性与并行性}

重复性任务在编程中很常见。

目前用得最广的并行范式是\textbf{单程序多数据}：代码只写一次，却以多个实例同时执行。除非显式同步，各实例彼此独立运行。Croqtile 采用这一范式，并用 \textbf{parallel-by} 块来抽象并行性。

有前后依赖的工作需要顺序循环，例如沿 K 切片累加部分结果。Croqtile 将 \co@with-in@ 声明的有界范围与 \co@foreach@ 执行的遍历分开，因此可以在遍历前后安排工作，多缓冲流水线就使用了这一结构。

循环处理必须依次执行的工作，\co@parallel-by@ 处理彼此独立的实例。


\co@parallel-by@ 块声明一组彼此独立的实例；空间说明符则规定这些实例如何分组。


\subsection*{并行执行块}

在 Croqtile 中，并行运行的代码放进\emph{并行执行块}，表示同一段代码有多个实例同时执行。





\subsection*{小结}

\co@parallel-by@ 块声明独立的实例，空间说明符把它们分组；掩码再决定其中哪些实例真正执行某条语句。

必须依次执行的工作仍然需要循环；2.7 节会介绍它们在 Croqtile 中的写法。





\subsection*{并行化与加掩码}

编写单程序多数据代码时，可能启动了固定数量的并行线程，但任务不能平均分配。掩码使有有效任务的线程执行，让其余线程保持非活跃状态。

在 Croqtile 代码里，掩码可以是隐式的，也可以是显式的。


\subsection*{显式掩码／条件}

线程掩码这个概念，从所有 SPMD 实例的整体视角描述执行状态。但在 SPMD 编程实践中，程序只有单线程视角，于是线程掩码表现为一个布尔表达式，通常是与线程标识符有关的比较结果。

后面的示例里，Croqtile 以实验性质支持对整个 \co@foreach@ 块加显式掩码，也就是把循环代码整段屏蔽。例如：

\begin{lstlisting}[language=croqtile,caption={显式掩码／条件}]
__co__  void foo() {
  parallel p by 6 {
    with q in [2] {
      foreach q if (p == 0) {
        // this iteration only work for parallel thread 0
      }
    }
  }
}
\end{lstlisting}


\subsection*{隐式掩码}

Croqtile 中有些代码结构自带隐式掩码，多层 \co@parallel-by@ 块就是典型。下面是一个例子：

\begin{lstlisting}[language=croqtile,caption={隐式掩码}]
__co__  void foo() {
  parallel p by 6 {
    // parallel-level-0:
    //   code here is as if implicitly masked 'if (q == 0)'
    parallel q by 2 {
      // parallel-level-1:
      //   code here is for all 'q' threads
    }
    // parallel-level-0:
    //   code here is as if implicitly masked 'if (q == 0)'
  }
}
\end{lstlisting}

本例有两层并行。位于 \co@parallel p by 6@ 块内、却在 \co@parallel q by 2@ 块之外的代码，只属于\emph{并行层次 0}。然而在某些目标平台上，实际会启动 2 个并行线程，分别执行\emph{并行层次 0} 或\emph{并行层次 1} 里的代码。于是，只属于\emph{并行层次 0} 的那部分，就好像被一个 C++ 的 \co@if (q == 0)@ 块保护着。

条件也可以用于保护整个 \co@foreach@ 循环，控制循环体是否执行。


\subsection*{实验特性：\texttt{within} 的可选 \emph{Where-Binding}}

还可以追加 \emph{where 子句}，在 \co@with-in@ 语句定义的多个有界变量之间施加 \textbf{Where-Binding} 约束，让不同的有界变量互为别名。例如：

\begin{lstlisting}[language=croqtile,caption={实验特性：within 的可选 Where-Binding}]
with {m, n} in [M, N], {n_p, k} in [N_P, K]
where n_p <-> n {
  // matmul implements with m, n, K. 'n_p' always refers to the same value as 'n'.
}
\end{lstlisting}

语法上，\co@<->@ 运算在两个有界变量之间建立 \emph{where-binding}。在 \co@with-in@ 块里，任何对 \co@n_p@ 的引用都指向 \co@n@，反过来也一样。


Croqtile 主要面向单程序多数据编程，但也通过线程掩码支持多程序多数据(MPMD)执行；发散代码就写在 \co@inthreads@ 块里。





\subsection*{SPMD 中的加掩码}

在 SPMD 编程中，常常需要屏蔽掉一部分\emph{并行线程}。比如，程序员把数据切成分块，交给多个线程并行或迭代处理，但可能除不尽，留下余量。

\begin{figure}[htbp]\centering
\includegraphics[width=\linewidth,height=0.56\textheight,keepaspectratio]{f34-masking-datachunk-dispatch.png}
\caption{13 个数据分块在 4 个并行线程上的分派}
\end{figure}

上图把 13 个数据分块分派给 4 个并行线程。要处理完全部分块需要 4 次迭代，但最后一次只有\emph{线程 0} 需要处理分块。前三次迭代的\textbf{线程掩码}都是 \co@1111@，只有最后一次是 \co@0001@（最低位对应线程 0）。线程掩码由若干布尔位组成，\co@1@ 表示\textbf{活跃}，\co@0@ 表示\textbf{非活跃}。所以本例中只把\emph{线程 0} 设为\emph{活跃}，避免在最后一次迭代中处理无效分块。


\subsection*{一致与发散}

加掩码是实现层面的概念。从 SPMD 编程的视角看，处理数据分块的代码通常被称为\textbf{发散}的，因为并非所有线程实例都执行同一段代码。程序员应当意识到，自己的 SPMD 代码由两部分组成：

\begin{itemize}[leftmargin=1.6em]
\item \textbf{一致代码}：\emph{所有线程}执行同一条代码路径。
\item \textbf{发散代码}：只有\emph{被掩码选中的线程}执行\emph{发散}代码路径。
\end{itemize}

现实中的不同需求，需要 SPMD 代码里有\emph{发散}代码来处理。换一种范式看，在 SPMD 程序里使用\emph{发散}代码，甚至能模拟多程序多数据的行为：MPMD 里不同线程执行不同代码，所以\emph{发散}在概念上把 SPMD 扩展到了 MPMD 的范围。

在 Croqtile 中，可以用 \co@inthreads@ 块来写发散代码。





\subsection*{多层次的加掩码}

线程掩码可以施加在不同的并行层次上。下面这个例子演示了这一点：

\begin{lstlisting}[language=croqtile,caption={多层次的加掩码}]
__co__  void foo() {
  parallel p by 3, q by 4 {
    inthreads (p < 2 && q == 0) {
      // divergent code for thread (0, 0), (1, 0)
      // thread-mask: 011-0001
    }
    inthreads (q == 1) {
      // divergent code for thread (0, 1), (1, 1), (2, 1)
      // thread-mask: 111-0010
    }
  }
}
\end{lstlisting}

这段代码用了两层并行：

\begin{itemize}[leftmargin=1.6em]
\item 外层（\co@p@）：并行数为 3
\item 内层（\co@q@）：并行数为 4
\end{itemize}

第一个 \co@inthreads@ 谓词（\co@p < 2 && q == 0@）把执行限制为：

\begin{itemize}[leftmargin=1.6em]
\item 外层并行 ID：\co@0, 1@（不含 2）
\item 内层并行 ID：仅 \co@0@
\end{itemize}

线程掩码 \co@011-0001@ 反映了这些约束，两个层次之间用连字符 \co@-@ 分隔。

第二个谓词（\co@q == 1@）不限制外层（\co@p@），允许所有外层 ID（\co@0, 1, 2@），只把内层限制为 \co@q == 1@。


\subsection*{小结}

加掩码让原本一致的并行区域内部也能出现发散执行。开发者既可以用 \co@inthreads@ 块做显式线程掩码，也可以借 \co@parallel-by@ 构造隐式控制，把执行限制在特定线程上。

加掩码确定哪些实例执行区域，其他实例是否在区域结束处等待，则取决于同步形式。第 3.8 节讨论异步区域：程序需要用事件表达数据依赖，不能依赖每个区域结束后的汇合。



\subsection{Worked Example: Tiled Matrix Multiplication}
\label{sec:3-4-1}

\subsection*{标量矩阵乘法（不使用 DMA）}

第一个版本用 \co@parallel@ 划分输出，用 \co@.at()@ 逐个读取元素，尚未加入显式的数据搬运。

\begin{lstlisting}[language=croqtile,caption={标量矩阵乘法（不使用 DMA）}]
__co__ s32 [128, 256] matmul(s32 [128, 256] lhs, s32 [256, 256] rhs) {
  s32 [lhs.span(0), rhs.span(1)] output;
  parallel {p, q} by [16, 64] {
    foreach {m, n} in [128 / #p, 256 / #q]
      output.at(p#m, q#n) = 0;
    foreach index = {m, n, k} in [128 / #p, 256 / #q, 256]
      output.at(p#m, q#n) += lhs.at(p#m, k) * rhs.at(k, q#n);
  }
  return output;
}
\end{lstlisting}

图~\ref{fig:f05-fig5_scalar_matmul} 画出了输出的分块网格。

\begin{figure}[htbp]\centering
\includegraphics[width=\linewidth,height=0.56\textheight,keepaspectratio]{f05-fig5_scalar_matmul_light.png}
\caption{[128, 256] 输出按 [16, 64] 网格划分，每块形状为 [8, 4]；图中只画出开头的一部分分块。}
\label{fig:f05-fig5_scalar_matmul}
\end{figure}

可以从输出形状开始，逐步核对并行划分、块内范围与全局坐标。

输出声明从 \co@lhs.span(0)@ 取得 128 行，从 \co@rhs.span(1)@ 取得 256 列。收缩维度出现在两个输入中，但不会成为输出维度。

声明 \co@parallel {p, q} by [16, 64]@ 为同一并行层次提供两个坐标轴。1,024 对索引将输出划分为每块八行四列。

循环将 \co@m@、\co@n@ 和 \co@k@ 绑定为 \co@index@ 的分量。前两者遍历输出块内位置，第三者遍历收缩维度。迭代次数直接来自输出划分：

\begin{itemize}[leftmargin=1.6em]
\item \co@128 / #p = 128 / 16 = 8@：每块的行数。
\item \co@256 / #q = 256 / 64 = 4@：每块的列数。
\item \co@256@：完整的收缩维度。
\end{itemize}

组合行索引 \co@p#m@ 的值为 \co@p * 8 + m@。p 遍历 16 个行分块，m 遍历每块八行，合起来覆盖第 0 到 127 行。列表达式 \co@q#n@ 使用相同规则，每块四列。

对每个输出位置，语句体沿 K 累加 \co@lhs.at(p#m, k) * rhs.at(k, q#n)@，所有贡献进入同一个输出元素。清零循环在同一并行实例中先执行，因此每个累加元素都由拥有它的实例清零。



\subsection{Tracing One Output Element}
\label{sec:3-4-2}

按照第 3.6.1 节 DMA 矩阵乘法的块与线程划分，追踪输出位置 \co@(37, 50)@。该例使用 8 × 16 的块网格，每块负责 16 × 16 的输出分块。这与前面标量示例中的 16 × 64 逻辑网格不同，但两者覆盖相同的输出形状。

各块把 128 行划分为 8 个大小为 16 的分块：\co@px = 37 / 16 = 2@，偏移 \co@qx = 37 % 16 = 5@。列方向把 256 划分为 16 个大小为 16 的分块：\co@py = 50 / 16 = 3@，偏移 \co@qy = 50 % 16 = 2@。于是块为 \co@(2, 3)@，线程为 \co@(5, 2)@。

对该线程，K 循环执行 16 次迭代。在 \co@tile_k = 0@ 这次迭代中，\co@dma.copy@ 把 \co@lhs@ 的第 32..47 行、K 的第 0..15 列装入共享的 \co@lhs_load@，把 \co@rhs@ 的 K 第 0..15 行、第 48..63 列装入共享的 \co@rhs_load@。内层 \co@foreach k in [16]@ 把 \co@lhs_load.data.at(5, k) * rhs_load.data.at(k, 2)@（k = 0..15）累加到 \co@output.at(37, 50)@。随后 \co@tile_k = 1@ 装入下一条 K 条带，依此类推。16 次迭代之后，\co@output.at(37, 50)@ 中就是完整的点积——与标量参考版本的答案相同。



\section{分块选择与数据表达式}
\label{sec:3-5}

\subsection*{\texttt{chunkat}、\texttt{subspan} 与 \texttt{view().from()}：三种选取分块的方式}

分块表达式需要回答两个问题：区域有多大，从哪里开始？使用 \co@chunkat@ 时，有界索引提供划分数量并选择一块；使用 \co@subspan(...).at(...)@ 时，分块形状显式给出，坐标按块计数；使用 \co@view(...).from(...)@ 时，形状和起始元素坐标都显式给出。窗口不落在规则分块边界上时，这种区别尤其重要。


\subsection*{\texttt{\#}：合成与范围}

一元形式 \co@#p@ 取得索引上界，二元形式 \co@p # i@ 将外层与内层索引组合起来，其值为 \co@p * #i + i@，上界为 \co@#p * #i@。若有八块、每块 16 个元素，第 2 块的块内第 5 个元素就对应 128 个元素范围中的第 37 个。


\subsection*{处理不规则访问}

尺寸能够整除且起点遵循固定网格时，规则分块很合适。其他核函数可能需要不同描述：专家批次从运行时偏移开始，卷积窗口相互重叠，最后一个分块的有效元素可能不足。扁平输入也可能需要先解释为多维视图，再参与计算。

下面的操作调整所选区域或形状，使外围算法仍保持原有结构。得到的视图仍需满足所选搬运操作的要求，包括形状、布局以及目标平台特定的对齐约束。


\subsection*{分块与流动}

\co@chunkat@ 表达式既可以作 \emph{DMA 语句}的源，也可以作目标。只把它用在源上，通常是把输入或缓冲区切成更小的块，搬到离处理器更近的存储；只用在目标上，则通常是搬运缓冲区，拼出输出的一部分。下面的代码是一个例子：

\begin{lstlisting}[language=croqtile,caption={分块与流动}]
__co__ auto foo(f32 [36] input) {
  f32 [36] output;
  parallel p by 6 : block {
    f0 = dma.copy input.chunkat(p) => shared;  // tile the source
    dma.copy f0.data => output.chunkat(p);     // un-tile the destination
  }
  return output;
}
\end{lstlisting}

本例把 \co@chunkat@ 用在目标，把 \co@shared@ 里的数据块搬出来，「反分块」地写进输出。注意元素类型必须一致，且源的任一维度都不能大于对应的目的维度；源比目的小是允许的。

同时在源和目标两端使用 \co@chunkat@ 表达式，目前几乎没有任何目标平台支持。


\subsection*{用 \texttt{\#} 合成两个有界变量以分块}

第 3.10.2 节 RGB 转灰度示例将工作实例索引 q 与顺序图像索引 n 组合起来。\co@q # n@ 为每个实例分配连续图像。N = 18、共有六个实例时，内层上界是 3，所以实例 1 处理图像 3、4、5。下面沿用这一划分。

共有六个工作实例时，q 的上界是 6，n 的上界是 N / 6。组合后的上界为 N，值为 \co@q * (N / 6) + n@。其余索引 h、w 选择空间分块，通道维保持完整。因此，实例分配与局部遍历共同确定源分块及对应输出位置。

交换操作数会改变分配方式。N = 18 时，随着 n 递增，\co@n # q@ 会让实例 1 处理图像 1、7、13，而 \co@q # n@ 选择的是连续的图像 3、4、5。运算符将这种调度选择直接写入数据表达式，实际性能取决于所得访问在目标平台上的表现。



\subsection{Arbitrary-Offset Windows: view and from}
\label{sec:3-5-1}

混合专家核函数中，不同专家经常处理不同数量的词元。每位专家数据的起点到运行时才确定，固定分块网格未必能描述所需切片。窗口可以先为切片确定自己的起点，再进行内部划分。

表达式 \co@view(M, N).from(row, col)@ 描述从给定元素坐标开始的 M × N 矩形。选择区域不必落在分块网格边界上，但使用该视图的搬运仍需满足自身的对齐和形状要求：

\begin{lstlisting}[language=croqtile,caption={常量偏移处的窗口}]
patch = matrix.view(16, 16).from(37, 50);
\end{lstlisting}

该表达式选取从第 37 行、第 50 列开始的 \co@[16, 16]@ 切片。原点可以是任意运行时值。

图~\ref{fig:f09-fig4_view_from} 展示了起点不在分块边界上的窗口。

\begin{figure}[htbp]\centering
\includegraphics[width=\linewidth,height=0.56\textheight,keepaspectratio]{f09-fig4_view_from_light.png}
\caption{图中用从 (2, 3) 开始的 [4, 4] 小窗口说明原理；正文代码从 (37, 50) 选取 [16, 16] 窗口时遵循相同规则。}
\label{fig:f09-fig4_view_from}
\end{figure}

\co@view@ / \co@from@ 的优势在于可组合性：选定窗口后，后面的所有操作——\co@subspan@、\co@chunkat@、\co@dma.copy@——都作用在这个窗口上，就像它是一个独立的张量。流水线保持不变，只是起始位置发生变化：

\begin{lstlisting}[language=croqtile,caption={运行时偏移处的窗口}]
expert_lhs = lhs.view(expert_M, K).from(expert_offset, 0);
dma.copy expert_lhs.subspan(WARP_M, TILE_K).at(block_m, iv_k) => shared;
\end{lstlisting}

在 MoE 栈里，每个专家的 token 批次都从动态行 \co@expert_offset@ 开始。\co@view@ / \co@from@ 只调整流水线入口的操作数；下面的 DMA、暂存和 MMA 循环，对每个专家都执行同一段代码。

划分数量已能描述所需区域时，使用 \co@chunkat@。起点是显式元素偏移时，使用 \co@view@ 和 \co@from@，再在窗口内部进行规则划分。这样就能将外层不规则选择与内部重复工作分开。



\subsection{Strided Tiles: subspan, step, and at}
\label{sec:3-5-2}

表达式 \co@subspan(M, K).at(i, j)@ 使用分块坐标。没有显式步长时，元素起点为 \co@(i * M, j * K)@。指定步长会改变相邻分块起点之间的距离，但不改变选中分块的大小。

加上 \co@.step(sM, sK)@ 会改写步长：分块之间相隔 \co@sM@ 行、\co@sK@ 列，可以按指定间隔排列：

\begin{lstlisting}[language=croqtile,caption={带步长的分块：.subspan、.step 与 .at}]
matrix.subspan(16, 16).step(32, 32).at(i, j);
\end{lstlisting}

图~\ref{fig:f09-fig5_subspan_step} 对比了紧凑排布与带步长的分块。

\begin{figure}[htbp]\centering
\includegraphics[width=\linewidth,height=0.56\textheight,keepaspectratio]{f09-fig5_subspan_step_light.png}
\caption{小图比较紧密排列的 [3, 3] 分块与 step(5, 5) 的间隔排列。正文的 [16, 16] 分块采用 step(32, 32)，留下 16 个元素的间隔。}
\label{fig:f09-fig5_subspan_step}
\end{figure}

分块 \co@(0,0)@ 从 \co@(0,0)@ 开始，而分块 \co@(1,0)@ 从第 32 行开始（不是第 16 行），留出 16 行的间隔。省略 \co@.step@ 时，步长等于分块尺寸，也就是紧凑排列的情形。

步长小于分块尺寸时，窗口相互重叠；步长更大时，窗口之间留下空隙。例如，16 行窗口每次前进八行，会与下一个窗口共享一半行。步长控制起点序列，不应与分块内部元素之间的存储步幅混淆。



\subsection{Layout Reinterpretation with span\_as}
\label{sec:3-5-4}

\subsection*{布局重新解释：\texttt{span\_as}}

有些核函数把数据当作扁平的一维条带载入，却要把它作为二维矩阵交给 MMA。朴素的做法是把数据复制到一个改过形状的缓冲区；\co@span_as@ 则不拷贝，直接就地重新解释现有缓冲区的形状：

\begin{lstlisting}[language=croqtile,caption={扁平条带被重新解释为矩阵}]
flat_buffer.span_as([rows, cols])
\end{lstlisting}

元素个数保持不变，只有逻辑秩改变，不搬运任何数据。

\begin{lstlisting}[language=croqtile,caption={使用重新解释结果的搬运语句}]
strip_load = dma.copy data.chunkat(tile) => shared;
tile_2d = strip_load.data.span_as([tile_m, tile_k]);
ma = mma.load tile_2d.chunkat(_, iv_warp);
\end{lstlisting}

\co@rows * cols@ 的乘积必须等于底层区域的元素数量。重解释形状会保留线性元素顺序，因此不能转置矩阵，也不能去掉行间填充。尺寸为常量时，编译器能在执行前检查元素数量；尺寸来自输入时，实际调用也必须满足同一约束。

图~\ref{fig:f09-fig7_span_as} 展示了这次重新解释。

\begin{figure}[htbp]\centering
\includegraphics[width=\linewidth,height=0.56\textheight,keepaspectratio]{f09-fig7_span_as_light.png}
\caption{16 个元素的缓冲区重解释为 [4, 4] 视图，线性顺序保持不变；图中将一维数据分成两行显示，便于阅读。}
\label{fig:f09-fig7_span_as}
\end{figure}



\subsection{Tile Selection: Summary and Syntax}
\label{sec:3-5-10}

\subsection*{小结}

\begin{table}[htbp]\centering\small
\begin{tabularx}{\linewidth}{@{}lXX@{}}
\toprule
\textbf{概念} & \textbf{语法} & \textbf{作用} \\
\midrule
软件 DMA & \texttt{dma.copy} / \texttt{dma.copy.swiz\textless{}N\textgreater{}} & 线程协作的分块传输 \\
重排 & 拷贝上的 \texttt{.swiz\textless{}N\textgreater{}} 与 \texttt{mma.load.swiz\textless{}N\textgreater{}} & 无布局冲突的共享内存布局；对不同搬运路径效果相同 \\
任意窗口 & \texttt{view(M,N).from(r,c)} & 不规整或运行时定位的切片 \\
带步长的分块 & \texttt{.subspan().step().at()} & 非紧凑布局、相互重叠的模板 \\
不完整分块 & \texttt{.zfill} & 把越界元素补零 \\
形状重新解释 & \texttt{span\_as([dims])} & 暂存缓冲区的零拷贝变形 \\
\bottomrule
\end{tabularx}
\end{table}


\subsection*{语法速览}

\begin{table}[htbp]\centering\small
\begin{tabularx}{\linewidth}{@{}lX@{}}
\toprule
\textbf{语法} & \textbf{含义} \\
\midrule
 &  \\
\texttt{dma.copy.swiz\textless{}N\textgreater{}} src =\textgreater{} dst & 采用以字节计的重排宽度 \texttt{N} 的拷贝 \\
\texttt{mma.load.swiz\textless{}N\textgreater{}} src & 与重排 \texttt{N} 一致的 MMA 操作数加载 \\
\texttt{.view(M, N).from(r, c)} & 任意偏移处的 \texttt{M x N} 窗口 \\
\texttt{.subspan(M, K).step(sM, sK).at(i, j)} & 带步长的分块选择 \\
\texttt{.zfill} & 把拷贝源中的越界元素补零 \\
\texttt{span\_as(dims)} & 把线性存储重新解释为带形状的张量 \\
\bottomrule
\end{tabularx}
\end{table}

3.6 节把上面的分块表达式用作搬运语句的源与目的。对于程序中确实需要纯 C++ 的部分，2.13 节汇集了 \co@__device__@ 函数、预处理守卫与注入的源代码文本。



\section{数据搬运：操作、配置与异步模式}
\label{sec:3-6}

\subsection*{到目前为止已具备的内容}

已介绍的构造可以拼成一副骨架：一个参数和返回值都带类型、带形状的 \co@__co__@ 函数；用 \co@parallel@ 和 \co@.at()@ 做逐元素计算；用 \co@lhs.span@ 从输入推导输出形状；以及主机侧用 \co@make_spandata@ / \co@.view()@ API 创建、传递张量。

2.9 节会引入 \co@dma.copy@ 和 \co@chunkat@，用来表达分块级的数据搬运。


\subsection*{数据搬运：分块搬运引擎}

逐元素写法使用标量索引：选一个位置、读输入、计算、写输出。

分块实现则先选择区域，将其暂存到计算使用的存储中，再在那里复用数据。Croqtile 直接表达这次搬运，也能将属于搬运过程的变换写入同一语句：

\begin{itemize}[leftmargin=1.6em]
\item \co@dma.copy@ 原样搬运一个分块。
\item \co@dma.transp<...>@ 在搬运的同时置换维度。
\item \co@dma.pad<...>@ 扩大分块并填充边缘。
\end{itemize}

这一模型的核心语法和规则，决定了 DMA 代码能否正确、高效地转换为目标代码。

\begin{figure}[htbp]\centering
\includegraphics[width=\linewidth,height=0.56\textheight,keepaspectratio]{f04-fig1_element_vs_block_light.png}
\caption{标量索引逐个读取元素；显式分块搬运先选择分块，再将其复制到 shared 存储。}
\end{figure}


\subsection*{其他操作：\texttt{transp} 与 \texttt{pad}}

普通复制保留源的形状和元素顺序。目标需要不同维度顺序时使用 \co@transp@，需要增加位置并填入指定值时使用 \co@pad@。目标形状也是操作的一部分，后续分块选择和计算应使用变换后的形状，而不是原始形状。


\subsection*{异步操作：DMA 语句}

除了并行执行，Croqtile 只允许一种固定形式的异步操作：DMA 语句。

概念上，DMA 语句与 SPMD 代码异步执行，行为很像 CPU 的异步线程，只是受 DMA 配置的约束（CPU 上的异步线程则可以执行一般程序代码）。

一条基本的 DMA 语句把源分块搬到目的存储，并返回一个异步句柄。

\begin{lstlisting}[language=croqtile,caption={一个异步 DMA 声明}]
global f32 [10] g_data;
local f32 [10] l_data;
f = dma.copy.async g_data => l_data;
// ... async operations
wait f;         // explicit wait
\end{lstlisting}

注意对 \co@f@ 的赋值。它是异步 DMA 实体的句柄，在 Croqtile 中称为该 DMA 操作的「期值」。程序员可以显式写出同步语句 \co@wait@，暂停当前线程直到期值到达；不写的话，线程就和 DMA 操作并行地继续执行。

显式声明临时缓冲区可能需要重复代码。Croqtile 提供了更简单的语法：

\begin{lstlisting}[language=croqtile,caption={带存储的同一种声明}]
global f32 [10] data;
f = dma.copy data => local;
... f.data;  // retrieve the 'local' data from the future
\end{lstlisting}

这里不必指定 DMA 把数据搬到哪个确切位置，只要说明目的内存类型。要取用已搬运的数据，调用期值的 \co@data@ 成员即可。另外，这个 dma 操作没有标 \co@.async@，所以对 \co@f@ 的 \co@wait@ 既不需要，也不能执行。

简写形式将分配交给编译器，但数据依赖并未改变。仍需通过期值确认目标，等待异步搬运完成，并检查后续操作是否接受所得形状与存储位置。下面的高级模式允许多次搬运同时进行，同时保留这些要求。


\subsection*{数据搬运语句}

Croqtile 要处理多个层次的数据搬运。它沿用了这个术语，并把数据搬运抽象为 \textbf{DMA 操作}。


\subsection*{操作类型}

Croqtile 在软件层把这些功能映射成不同的 \textbf{DMA 操作类型}。支持的操作包括：

\begin{itemize}[leftmargin=1.6em]
\item \co@dma.copy@：直接拷贝扁平内存。
\item \co@dma.pad@：搬运带形状数据的同时做填充。
\item \co@dma.transp@：搬运带形状数据的同时做转置。
\end{itemize}

填充、转置这些细节，作为 DMA 操作的参数来配置。例如：

\begin{lstlisting}[language=croqtile,caption={操作类型}]
global f32 [32, 16, 9] input;
dma.transp<0, 2, 1> input => shared;  // Result shape [32, 9, 16]
dma.pad<{1, 0, 3}, {0, 1, 2}, {0, 0, 0}, 0.1f> input => shared; // Result shape [33, 17, 14]
\end{lstlisting}

这里的 \textbf{DMA 配置}用 \co@<...>@ 括起，内容随操作不同而变化。以某个目标平台为例，具体的配置语法和限制如下：

要注意，\co@dma.copy@ 之外的 DMA 类型，支持程度因平台而异。比如在某些平台上，它可能映射到专用搬运引擎，也可能是对多个线程加载指令的编排。


数据搬运对 \emph{tileflow 程序}至关重要。\emph{DMA 语句}为各个层次各提供一种形式，从单次拷贝一直到链式非阻塞传输。











\subsection*{完全非阻塞 DMA 模式（Croqtile 中的链模式）}

使用前一次搬运结果的传输，必须在前一次完成后才能开始。\co@after@ 子句记录这项依赖，而不要求发起操作的代码当场等待。下面的例子将暂存缓冲区的复制与此前操作连接起来：

\begin{lstlisting}[language=croqtile,caption={完全非阻塞 DMA 模式（Croqtile 中的链模式）}]
out_store = dma.copy.async l2_out => output.chunkat(m_tile, n_tile) after out_store_s;
\end{lstlisting}

这个例子里：

\begin{itemize}[leftmargin=1.6em]
\item \co@out_store@ 是把数据从 \co@l2_out@ 搬到 \co@output.chunkat(m_tile, n_tile)@ 的异步 DMA 操作。
\item \co@out_store_s@ 是另一项 DMA 操作或事件，必须在 \co@out_store@ 开始前完成。
\item \co@after out_store_s@ 语法规定，\co@out_store@ 只能在 \co@out_store_s@ 完成后启动，从而保证主线程不被阻塞。
\end{itemize}

依赖尚未满足时，发起操作的代码可以继续执行，但目标并不会立即变得可读，源也不能因此提前复用。后续消费者仍需要等最终搬运完成。阅读搬运链时，应沿源到目标检查每次交接，包括最后回到普通计算的那一次。












\subsection{Worked Example: DMA Matmul}
\label{sec:3-6-1}

\subsection*{DMA 矩阵乘法：共享存储中的分块}

\begin{lstlisting}[language=croqtile,caption={DMA 矩阵乘法：共享存储中的分块}]
__co__ s32 [128, 256] matmul(s32 [128, 256] lhs, s32 [256, 256] rhs) {
  s32 [lhs.span(0), rhs.span(1)] output;
  parallel {px, py} by [8, 16] : block {
    parallel {qx, qy} by [16, 16] : thread
      output.at(px#qx, py#qy) = 0;
    foreach {tile_k} in [16] {
      lhs_load = dma.copy lhs.chunkat(px, tile_k) => shared;
      rhs_load = dma.copy rhs.chunkat(tile_k, py) => shared;
      parallel {qx, qy} by [16, 16] : thread {
        foreach k in [256 / #tile_k]
          output.at(px#qx, py#qy) += lhs_load.data.at(qx, k) * rhs_load.data.at(k, qy);
      }
    }
  }
  return output;
}
\end{lstlisting}

\begin{figure}[htbp]\centering
\includegraphics[width=\linewidth,height=0.56\textheight,keepaspectratio]{f05-fig6_dma_matmul_light.png}
\caption{每个 block 为当前 K 切片装入两个 [16, 16] 操作数分块，再由 256 个线程累加该 block 的 [16, 16] 输出分块。}
\end{figure}

标量版本一次索引一个元素；分块版本先把一个分块复制到共享存储，再在这块数据上计算：

\begin{itemize}[leftmargin=1.6em]
\item 外层 \co@parallel@ 现在用 \co@: block@ 和花括号形式的索引 \co@{px, py}@。8 × 16 的块网格覆盖整个输出。
\item 对 \co@tile_k@ 的 \co@foreach@ 分 16 步走完 K 维。每一步用 \co@dma.copy ... => shared@ 把 \co@lhs@ 和 \co@rhs@ 的一条条带复制到 \co@shared@ 存储。
\item 内层 \co@parallel {qx, qy} by [16, 16] : thread@ 在每个块里创建 256 个线程，每个线程负责块内分块的一个输出元素。
\item 算术从 \co@lhs_load.data@ 和 \co@rhs_load.data@ 读取，也就是共享内存里的副本，而不是全局的 \co@lhs@ 和 \co@rhs@。
\end{itemize}

合成索引分两层：\co@px#qx@ 把块索引 \co@px@ 和线程索引 \co@qx@ 合成为全局行，\co@py#qy@ 对列做同样的事。

块网格 \co@[8, 16]@ 为每块分配 16 行、16 列输出。块内 \co@[16, 16]@ 个线程各负责一个输出元素。K 方向共有 16 个分块，每块含 16 个元素，共同覆盖大小为 256 的完整收缩维度。

数据复用发生在块内：装入的左侧分块为多个输出列供数，右侧分块为多个输出行供数。K 循环遍历各切片一次，将贡献累加到同一组输出位置。与标量示意代码一样，输出在 K 循环之前清零。暂存改变了内层循环的读取位置，点积本身保持不变。



\subsection{Multi-Buffering: dma.any, swap, and rotate}
\label{sec:3-6-2}

\subsection*{用 \texttt{swap} 做双缓冲}

给 K 循环准备两个逻辑缓冲区：计算使用缓冲区 0 时，DMA 将下一个分块装入缓冲区 1。计算完成后交换句柄，原来的“下一个”变成“当前”，释放的槽位供下一次装入使用。

在 Croqtile 中，这一过程由 \co@dma.copy.async@（非阻塞拷贝）、\co@dma.any@（占位期值）、\co@swap@（交换期值句柄）和一个三阶段循环来表达：

\begin{lstlisting}[language=croqtile,caption={用 swap 做双缓冲}]
__co__ auto matmul(s32 [M, K] lhs, s32 [K, N] rhs) {
  s32 [lhs.span(0), rhs.span(1)] output;
  parallel {px, py} by [8, 16] : block
    parallel {qx, qy} by [16, 16] : thread {
    output.at(px#qx, py#qy) = 0;
    with tile_k in 16 {
      // Prologue: start loading tile 0
      lf0 = dma.copy.async lhs.chunkat(px, tile_k) => shared;
      rf0 = dma.copy.async rhs.chunkat(tile_k, py) => shared;
      // Placeholder futures for buffer 1
      lf1 = dma.any;
      rf1 = dma.any;
      // Steady state: load next tile while computing on current
      foreach tile_k(1:) {
        lf1 = dma.copy.async lhs.chunkat(px, tile_k) => shared;
        rf1 = dma.copy.async rhs.chunkat(tile_k, py) => shared;
        wait lf0, rf0;
        foreach k in [256 / #tile_k]
          output.at(px#qx, py#qy) += lf0.data.at(qx, k) * rf0.data.at(k, qy);
        swap(lf0, lf1);
        swap(rf0, rf1);
      }
      // Epilogue: compute on the last loaded tile
      foreach k in [256 / #tile_k]
        output.at(px#qx, py#qy) += lf0.data.at(qx, k) * rf0.data.at(k, qy);
    }
  }
  return output;
}
\end{lstlisting}


\subsection*{三个阶段}

前序阶段发起第 0 块的装入，用 \co@lf0@ 和 \co@rf0@ 表示。此时还没有上一块可供计算，这两次搬运为后面的循环建立初始状态。

稳定阶段先用 \co@lf1@ 和 \co@rf1@ 发起下一块的装入，再等待当前块就绪，并从 \co@lf0@ 和 \co@rf0@ 读取计算。两对期值对应不同缓冲区，因此下一块的搬运可以与当前计算重叠，并不要求下一块先装入完成。当前计算结束后，\co@swap@ 将下一对变为下一次迭代的当前对。

循环停止发起搬运后，收尾阶段处理最后装入的分块。读取之前，仍需确认它已装入完成。上面的清单用于展示期值轮换，省略了最后一次等待和输出初始化；实际执行累加时必须补齐这两项，也必须像稳定阶段一样保证读取时数据已就绪。


\subsection*{\texttt{swap}：交换的是名字，不是字节}

操作 \co@swap(lf0, lf1)@ 交换表示当前和下一缓冲区的句柄，不复制元素，也不等待搬运。三个缓冲区可用 \co@rotate(f0, f1, f2)@ 轮换句柄。无论哪种情况，都应在交换后继续追踪句柄，确定读取之前需要等待哪次搬运。


\subsection*{\texttt{dma.any}：占位期值}

\co@dma.any@ 在尚未发起搬运时先声明期值名字，使两个句柄都能位于包含循环与收尾代码的作用域中。占位期值并不包含已就绪的分块，必须先用实际搬运替换它，才能使用其数据。


\subsection*{用 Croqtile 做多缓冲优化}

要激进地优化高性能计算核函数，一个有效的办法是让计算与 DMA（直接内存访问）数据搬运重叠。这种做法叫\textbf{多缓冲优化}，通常要求程序员管理多个缓冲区。

在 Croqtile 中，这种做法需要两个关键原语：占位期值 \textbf{dma.any}，以及 \textbf{swap/rotate} 函数：

\begin{lstlisting}[language=croqtile,caption={用 Croqtile 做多缓冲优化}]
__co__ s32 [6, 17, 128] ele_add(s32 [6, 17, 128] lhs,
                                s32 [6, 17, 128] rhs) {
  s32[lhs.span] output;
  parallel p by 6 {
    with index = {x, y} in [17, 4] {
      local s32 [rhs.span / {#p, #x, #y}] l_out;
      foreach x { // first bunch
        lfA = dma.copy.async lhs.chunkat(p, x, y) => local;
        rfA = dma.copy.async rhs.chunkat(p, x, y) => local;
        lfB = dma.any;  // dummy DMA used for multi-buffering
        rfB = dma.any;
        foreach y(1:) {
          lfB = dma.copy.async lhs.chunkat(p, x, y) => local;
          rfB = dma.copy.async rhs.chunkat(p, x, y) => local;
          wait lfA, rfA;  // wait another bunch
          call kernel(lfA.data, rfA.data, l_out, |lfB.span|);
          dma.copy l_out => output.chunkat(p, x, y - 1);
          swap(lfA, lfB); // exchange futures
          swap(rfA, rfB);
        }
        lf = select(#y % 2, lfB, lfA);
        rf = select(#y % 2, rfB, rfA);
        wait lf, rf;  // handle the last bunch
        call kernel(lf.data, rf.data, l_out, |lfB.span|);
        dma.copy l_out => output.chunkat(p, x, y(-1));
      }
    }
  }
  return output;
}
\end{lstlisting}

这个加法核函数为每份输入维护两组局部缓冲区。A 句柄表示当前操作数，B 句柄接收下一次装入。当前结果计算并写回后，通过交换句柄推进两组角色。左右输入沿同一调度前进，使每次加法使用相互对应的分块。

实现把数据搬运和计算分成三个阶段：

\begin{itemize}[leftmargin=1.6em]
\item \textbf{前奏}（第 8--11 行）：先（预）装入第一块数据。本例中缓冲区 A 被预先装好。
\item \textbf{主体}（第 12--20 行）：处理上一轮迭代或前奏预装的数据，同时预装下一轮要用的数据。
\item \textbf{尾声}（第 21--25 行）：只处理最后一块数据。
\end{itemize}

占位期值 \co@lfB@ 和 \co@rfB@ 在内层循环前声明，因此收尾代码仍能引用它们。循环先用实际搬运结果替换占位值，再使用其数据。若只在循环内部声明这些名字，处理最后一块时它们就已超出作用域。

第 18--19 行的 \textbf{swap} 语句交换期值，让第 15 行的 \textbf{wait} 语句能正常工作。注意，只有对应相同 DMA 操作的期值才能交换。另外，第 12 行 \emph{foreach} 的迭代范围也调过：\co@foreach y(1:)@ 表示对\emph{有界变量} \co@y@ 的循环从第二个取值开始，一直到上界。第 17 行的 \emph{chunkat} 表达式同样调过：\co@y - 1@ 取 \co@y@ 的当前值减 1，让主体阶段和前奏阶段正确衔接。这个调整是必要的：迭代次数要减 1，而预装载操作应该在计算「当前」数据块时装入「下一」块。也可以改用 \co@foreach y(:-1)@，再相应调整第 13--14 行预装载语句的 chunkat 表达式，效果一样。

第 25 行的 \co@y(-1)@ 取出 \co@y@ 在界内的「最后一个取值」。

第 21--22 行的 \textbf{select} 表达式处理\emph{有界变量} \co@y@ 上界为奇数和偶数两种情况，无论哪种都能取得最后一块数据对应的期值。

如果想用两个以上的缓冲区，就要用 \textbf{rotate} 语句代替 \textbf{swap}。



\section{收缩与布局约定}
\label{sec:3-7}

\subsection*{张量收缩：\texttt{mma} 语法}

3.4 节的分块矩阵乘法是按 CPU 的方式计算内部收缩的：一个 \co@foreach k@ 循环，每次迭代把两个标量相乘，累加到同一个输出元素上。这种写法可以得到正确结果。

真正的任务是\textbf{二维张量收缩}——C += A × B，其中 A、B、C 都是形状固定的小矩阵分块。它是每个 GEMM、每个以 im2col 矩阵乘表达的卷积、每个注意力头 QKT 的内层核函数。

困难在于矩阵元素如何分布到参与协作的操作数片段中。布局声明必须与数据布局一致，例如把列主序分块当作行主序片段使用，就可能得到错误结果。下面的颜色只示意分布，不表示某条指令的精确映射。

图~\ref{fig:f06-fig2_register_loading} 展示了操作数分块如何分布到协作的各个 lane 上。

\begin{figure}[htbp]\centering
\includegraphics[width=\linewidth,height=0.56\textheight,keepaspectratio]{f06-fig2_register_loading_light.png}
\caption{操作数片段分布在协作的线程之间。颜色仅示意这种分布，不代表某条指令的精确寄存器映射。}
\label{fig:f06-fig2_register_loading}
\end{figure}

Croqtile 用 \co@fill@、\co@load@、乘加和 \co@store@ 表达收缩，由编译器处理片段元素在所选协作范围内的分布。程序员仍需选择兼容的分块形状并声明操作数布局，这层抽象以这些选择正确为前提。

图~\ref{fig:f06-fig3_mma_syntax} 展示了这四个步骤以及它们沿 K 维的重复。

\begin{figure}[htbp]\centering
\includegraphics[width=\linewidth,height=0.56\textheight,keepaspectratio]{f06-fig3_mma_syntax_light.png}
\caption{先初始化累加器，再装入两个操作数、执行累加，最后写回；沿 K 切片重复装入和累加步骤。}
\label{fig:f06-fig3_mma_syntax}
\end{figure}


\subsection*{扩大协作范围}

无论多少线程协作完成一次收缩，四步语法都不变。变的是\textbf{协作范围}——以下示例围绕这一变化展开。


\subsection*{不完整分块、守卫与 \texttt{.zfill}}

真实核函数常常以不完整的分块收尾。关键区别在于\textbf{源分块形状}和\textbf{目的分块形状}是否相同。

例如：

\begin{lstlisting}[language=croqtile,caption={不完整分块、守卫与 .zfill}]
dma.copy.swiz<32>.zfill
  lhs.view(ROWS_M, TILE_K).from(base_m, base_k)
    => lhs_load_s;
\end{lstlisting}

这里 \co@lhs_load_s@ 的形状可能是 \co@[group_M, TILE_K]@，源视图的形状却是 \co@[ROWS_M, TILE_K]@；在最后一个分块上，\co@ROWS_M@ 可能小于 \co@group_M@。这是源区域与目的区域之间的形状差异，此时 \co@.zfill@ 的意思是：拷贝有效的源行，把目的分块剩下的行写零。

这正是 \co@.zfill@ 的正确用法。

形状相同时，情况是这样：

\begin{itemize}[leftmargin=1.6em]
\item 源视图和目的分块形状相同时，编译器已经自动给运行时尾部加了守卫。
\item 这时它会生成掩码，让越界元素既不被读，也不被写。
\item 对把 \co@[100]@ 张量切成 \co@[64]@ 分块这类情况，这已经够了：最后一个分块只有 36 个有效元素。
\end{itemize}

\co@.zfill@ \textbf{不是}处理「任意不完整分块」的通用机制。它专门针对目的分块大于有效源区域、且目的端多出的部分必须填零的情形。

\co@.zfill@ 为什么重要：

\begin{itemize}[leftmargin=1.6em]
\item 消费者只读有效区域时，编译器的自动掩码通常就够了。
\item 消费者要读整个目的分块时，尤其是 MMA 片元，无效区域里必须是已知值。
\item \co@.zfill@ 把多出来的那块区域置零。
\end{itemize}

当有效源区域小于目的分块，且消费者会把目的端当作完整的稠密分块时，应使用 \co@.zfill@。这是带动态行数或带填充边缘分块的 GEMM 式流水线里的常见情形。

源和目的形状本来就相同时再加 \co@.zfill@，编译器可能警告说这个 \co@.zfill@ 是多余的；反过来，形状不同时，编译器也可能警告需要 \co@.zfill@，免得未覆盖区域里出现未定义值。


\subsection*{\texttt{.swiz\textless{}N\textgreater{}}：面向消费者的共享内存布局}

典型的 GEMM 暂存代码如下：

\begin{lstlisting}[language=croqtile,caption={.swiz N：面向消费者的共享内存布局}]
dma.copy.swiz<128> lhs.subspan(WARP_M, TILE_K).at(block_m, iv_k) => lhs_load_s;
ma = mma.load.swiz<128> lhs_load_s.chunkat(_, iv_warp);
\end{lstlisting}

装载侧和计算侧的重排必须一致。用 \co@.swiz<128>@ 暂存的分块，消费者通常也要用匹配的重排来加载。

\co@.swiz<N>@ 的取值决定共享内存中分块元素的排列，必须与 \co@mma.load.swiz<N>@ 一致。


\subsection*{重排与布局冲突}

交错布局改变暂存分块的元素在共享存储中的排列。逻辑矩阵形状不变，但读取方必须使用匹配的地址映射。Croqtile 的拷贝操作都提供这一选择：

\begin{lstlisting}[language=croqtile,caption={一次带重排的拷贝}]
dma.copy.swiz<128> src => dst;       // software DMA with swizzle
dma.copy.swiz<128> src => dst;
\end{lstlisting}

MMA 的读取路径必须使用相同的 \co@swiz<N>@，让地址与暂存的布局一致：

\begin{lstlisting}[language=croqtile,caption={一次带重排的 mma 加载}]
ma = mma.load.swiz<128> lhs_load_s.chunkat(_, iv_warp);
\end{lstlisting}

重排是拷贝自身的性质，而不是某个引擎独有的特性：\co@dma.copy.swiz<N>@ 产生的共享内存布局与数据如何到达无关。

模板参数选择目标支持的交错模式。示例使用编号模式，其中模式 0 表示恒等映射。有效模式及兼容的分块尺寸取决于目标平台，搬运与 MMA 装入应采用同一套已记录的约定。

复制与 \co@mma.load.swiz<N>@ 必须对布局达成一致。用不兼容模式读取交错分块，会把相同字节解释为不同矩阵位置，导致数值错误。修改任一操作时，都应将两者匹配视为正确性要求。

图~\ref{fig:f09-fig2_swizzle} 说明了对布局做这种改变为什么重要。

\begin{figure}[htbp]\centering
\includegraphics[width=\linewidth,height=0.56\textheight,keepaspectratio]{f09-fig2_swizzle_light.png}
\caption{八个 bank 的简化示意：同一 bank 的不同地址访问需要串行处理，重映射则分散访问。复制与 MMA 装入必须采用一致的 swizzle。}
\label{fig:f09-fig2_swizzle}
\end{figure}


\subsection*{稠密 FP16 之外：四步还能表达什么}

上面的例子对稠密 FP16 分块使用 \co@mma.row.row@。同样的四步模式可以扩展到基本形式覆盖不到的工作负载。

稀疏收缩还需要一个\textbf{元数据操作数} \co@me@，用来标明哪些元素非零：

\begin{lstlisting}[language=croqtile,caption={一次稀疏收缩}]
mma.row.row.sp mc, ma, mb, me;
\end{lstlisting}

\co@.sp@ 后缀加上元数据操作数，其余仍是填充、加载、相乘、存储的顺序。任何布局组合都能用，比如 \co@mma.row.col.sp@。\co@me@ 和 A、B 一起，从单独的元数据张量里载入。

对带分块缩放值的量化操作数，收缩还需要所选 A、B 分块对应的缩放值。\co@.scale@ 形式将缩放与矩阵操作一起表达：

\begin{lstlisting}[language=croqtile,caption={一次带缩放的收缩}]
mma.row.row.scale mc, ma, mb, sc_a, sc_b;
\end{lstlisting}

收缩之后，每个结果元素由 \co@sc_a@ 和 \co@sc_b@ 缩放——不再需要单独的反量化核函数。如果缩放来源和标准的融合路径不同，缩放也可以写成\textbf{独立语句}：

\begin{lstlisting}[language=croqtile,caption={另一种收缩布局}]
mma.row.row mc, ma, mb;
mma.scale mc, sc_a, sc_b;
\end{lstlisting}

独立的 \co@mma.scale@ 出现在某些 MoE 与混合精度核函数中。

\co@dma.copy.swiz<N>@ 和 \co@mma.load.swiz<N>@ 里的 \co@N@ 必须一致，不一致就会读出错误数据。输出方面，\co@mma.store.transp mc, dst@ 会把累加器的行和列互换后写出，适合下一阶段需要列主序数据的情形。

所有变体汇总在下表中，方便查阅。



\subsection{Worked Example: One Cooperative Group}
\label{sec:3-7-1}

\subsection*{一组一分块：最简单的 MMA 矩阵乘法}

最小示例将一个输出分块分配给用 \co@: group@ 表示的协作组。函数接收左右操作数和调用方提供的输出缓冲区。沿着一个输出分块阅读代码：初始化累加器，遍历各 K 切片，将累加分块存入共享存储，最后复制到对应输出位置。

\begin{lstlisting}[language=croqtile,caption={一组一分块：最简单的 MMA 矩阵乘法}]
#define TILE_M 16
#define TILE_N 16
#define TILE_K 16
#define WARP_M 16
#define WARP_N 16
#define WARP_K 16
#define MMA_M 16
#define MMA_N 16
#define MMA_K 16
__co__ void matmul_wmma(global f16 [M, K] lhs, global f16 [N, K] rhs, global f16 [M, N] output) {
  parallel {block_m, block_n} by [cdiv(M, TILE_M), cdiv(N, TILE_N)] : block {
    shared f16 [TILE_M, TILE_N] output_s;
    parallel {warp_m, warp_n} by [cdiv(TILE_M, MMA_M), cdiv(TILE_N, MMA_N)] : group {
      mc = mma.fill 0.0;
      foreach {iv_k} in [cdiv(K, TILE_K)] {
        lhs_load_s = dma.copy lhs.subspan(TILE_M, TILE_K).at(block_m, iv_k) => shared;
        rhs_load_s = dma.copy rhs.subspan(TILE_N, TILE_K).at(block_n, iv_k) => shared;
        foreach iv_warp_k in [cdiv(TILE_K, MMA_K)] {
          ma = mma.load lhs_load_s.chunkat(warp_m, iv_warp_k);
          mb = mma.load rhs_load_s.chunkat(warp_n, iv_warp_k);
          mma.row.row mc, ma, mb;
        }
      }
      mma.store mc, output_s.subspan(MMA_M, MMA_N).at(warp_m, warp_n);
    }
    dma.copy output_s => output.subspan(TILE_M, TILE_N).at(block_m, block_n);
  }
}
\end{lstlisting}

函数返回 \co@void@，因为它直接写入参数 \co@output@。调用方提供该缓冲区，并在操作完成后读取结果。函数内部的 \co@output_s@ 是 \co@mma.store@ 与最终复制之间使用的暂存缓冲区。

向上取整除法 \co@cdiv(M, TILE_M)@ 计算 M 方向分块数量，包括可能不足一整块的尾部。\co@(block_m, block_n)@ 标识当前块负责的输出分块。操作数选择复用这些输出坐标，并加入当前 K 切片坐标。

当矩阵尺寸与分块尺寸都为 16 时，块网格为 1 × 1。更大的矩阵会产生更多输出分块。编译这段清单前，需要补充 \co@MMA_M@、\co@MMA_N@ 和 \co@MMA_K@ 的定义，并使它们与所选片段尺寸一致。

3.5 节详细讨论 \co@subspan@ 和 \co@chunkat@ 的区别。

\begin{keypoint}
\co@subspan(h, w)@ 声明固定大小的矩形视图形状，\co@.at(i, j)@ 从父张量中选出想要的视图实例。换言之，\co@lhs.subspan(TILE_M, TILE_K).at(block_m, iv_k)@ 的意思是「取位于 \co@(block_m, iv_k)@ 的 \co@[TILE_M, TILE_K]@ 分块」。需要显式分块形状时用 \co@subspan@；需要按已有维度的等分来切时用 \co@chunkat@。
\end{keypoint}

累加器跨越整个 K 循环。若在循环内初始化，就会丢弃此前的贡献；若在循环内写回，得到的只是部分结果。将初始化放在循环前、写回放在循环后，能清楚表示输出分块的累加过程。



\subsection{MMA Syntax Summary}
\label{sec:3-7-9}

\subsection*{语法速览}

\begin{table}[htbp]\centering\small
\begin{tabularx}{\linewidth}{@{}lX@{}}
\toprule
\textbf{语法} & \textbf{含义} \\
\midrule
\texttt{mc = mma.fill 0.0} & 把累加器分块初始化为零 \\
\texttt{mma.fill.f16 0.0f} / \texttt{mma.fill.f32 0.0f} & 显式指定精度的累加器 \\
\texttt{mma.load.swiz\textless{}N\textgreater{}} src & 按重排模式加载 \\
\texttt{mma.row.row mc, ma, mb} & C += A × B（均为行主序） \\
\texttt{mma.row.col mc, ma, mb} & C += A × B（A 行主序，B 列主序） \\
\texttt{mma.row.row.sp mc, ma, mb, me} & 带元数据操作数的稀疏 MMA \\
\texttt{mma.row.row.scale mc, ma, mb, sc\_a, sc\_b} & 融合的 MMA 与逐分块反量化 \\
\texttt{mma.scale mc, sc\_a, sc\_b} & 独立的 MMA 后置缩放 \\
\texttt{mma.store.transp mc, dst} & 转置后写出累加器 \\
\texttt{cdiv(a, b)} & 向上取整除法 \\
\texttt{\_\_co\_\_ void fn(...)} & 就地写出结果的核函数 \\
\bottomrule
\end{tabularx}
\end{table}

阅读修饰符时，应同时追踪数据路径。交错装入必须匹配复制产生的布局，稀疏操作需要同一操作数分块的元数据，缩放值也必须对应正在累加的分块。语句写得简短，并不意味着可以忽略这些关系。



\subsection{MMA Design Summary}
\label{sec:3-7-10}

\subsection*{小结}

\begin{table}[htbp]\centering\small
\begin{tabularx}{\linewidth}{@{}lX@{}}
\toprule
| \textbf{主题} & \textbf{要点} | \\
\midrule
| --- & --- | \\
| 二维张量收缩 & 在分块状操作数上做 C += A × B —— 通用的内层核函数 | \\
| 四步抽象 & \textbf{填充 → 加载 → 相乘 → 存储}；编译器为各目标平台处理片元布局 | \\
| 扩大协作规模 & 语法保持不变 | \\
| 布局约定 & \co@mma.row.row@、\co@mma.row.col@ 等 | \\
| 稀疏与量化 & \co@.sp@ 增加元数据操作数；\co@.scale@ 融合逐分块反量化 | \\
| 重排与流水线 & \co@mma.load.swiz<N>@ 匹配带重排的共享布局 | \\
\bottomrule
\end{tabularx}
\end{table}

收缩执行得更快，但装载和计算仍然轮流进行——3.8 节会引入\textbf{角色专门化}，让不同的线程组各司其职：一组装载数据，另一组做计算，使内存访问和数学运算在时间上重叠。



\section{协调、流水线与持久调度}
\label{sec:3-8}

\subsection*{控制流：发散、并发与协同}

每一种面向并行机的语言都必须回答同一个问题：\textbf{当不同线程需要做不同的事情时，会发生什么？}

CPU 的各个核心具有独立的程序计数器；共享指令流的线程遇到发散路径时，则需要依次执行各条路径，并屏蔽当前不执行的线程。

所有线程都求值该条件，条件为假的线程被屏蔽（停用）；这就是\textbf{谓词执行}——发散路径可能需要额外执行时间。

分块矩阵乘法还包含职责不同的工作：一个组装入下一批操作数，另一个组使用当前操作数计算。两组需要在明确的交接点之间独立推进。运行时条件描述实例选择哪个分支，而流水线还需要表达这些独立职责及其依赖关系。

Croqtile 引入两个额外的控制流原语来填补这个空白：

\begin{itemize}[leftmargin=1.6em]
\item \textbf{\texttt{inthreads.async}} —— \textbf{结构化并发区域}：在编译时把线程分成若干组，各组\textbf{同时运行不同的程序}，编译器为它们生成彼此独立的指令流；
\item \textbf{\texttt{shared event} / \texttt{wait} / \texttt{trigger}} —— \textbf{区域间通信}：轻量级的同步令牌，让并发区域能安全地通信。
\end{itemize}


\subsection*{用 \texttt{inthreads.async} 构造结构化并发区域}

形式 \co@inthreads.async@ 将不同线程组分配到各自推进的区域，谓词标识负责该区域的组。生产者和消费者的职责固定，而进度取决于哪些缓冲区就绪，适合用这种方式表达。

\begin{lstlisting}[language=croqtile,caption={用 inthreads.async 构造结构化并发区域}]
parallel p1 by 2 : thread {
  inthreads.async (p1 == 0) {
  }
  inthreads.async (p1 == 1) {
  }
}
\end{lstlisting}

与 \co@if@ 的区别是结构性的，不只是性能上的：

\begin{table}[htbp]\centering\small
\begin{tabularx}{\linewidth}{@{}lXX@{}}
\toprule
 & \textbf{\texttt{if}（谓词执行）} & \textbf{\texttt{inthreads.async}}（结构化并发） \\
\midrule
\textbf{判定时机} & 运行时——每个线程都求值条件 & 编译时——线程归属已固定 \\
\textbf{指令流} & 一段程序；发散线程被屏蔽 & 每个区域各有一段程序 \\
\textbf{编程语言类比} & 任意语言中的 \texttt{if}/\texttt{else} & 结构化并发中的 \texttt{async}/\texttt{spawn}（Trio、Go 协程、Cilk） \\
\bottomrule
\end{tabularx}
\end{table}

区域由词法作用域限定，编译器可以从程序结构中识别参与者。代码不会在运行时创建数量不定的任务，每个区域都在已声明的并行范围内执行。

不加 \co@.async@ 时，区域在继续后续代码之前与外围执行汇合。加上 \co@.async@ 后，其他区域可以独立推进，它们之间的数据交接需要显式同步协议。

下图展示了实际效果——生产者的 DMA 和消费者的 MMA 在时间上重叠：

图~\ref{fig:f07-fig1_role_comparison} 对比了这两种调度。

\begin{figure}[htbp]\centering
\includegraphics[width=\linewidth,height=0.56\textheight,keepaspectratio]{f07-fig1_role_comparison_light.png}
\caption{两种调度都包含四组装入与计算任务。并发的生产者和消费者在遵守依赖关系的前提下重叠执行。}
\label{fig:f07-fig1_role_comparison}
\end{figure}


\subsection*{用事件做区域间通信}

\co@inthreads.async@ 创建并发区域之后，这些区域需要一种通信方式。Croqtile 提供了\textbf{事件}。

\begin{lstlisting}[language=croqtile,caption={用事件做区域间通信}]
shared event full;
shared event empty;
\end{lstlisting}

事件有两种操作：

\begin{itemize}[leftmargin=1.6em]
\item \textbf{\texttt{trigger name}} —— 发出条件已满足的信号（例如「数据已就绪」）
\item \textbf{\texttt{wait name}} —— 阻塞直到相应的 \co@trigger@ 触发
\end{itemize}

生产者写完一个分块后调用 \co@trigger full@，发出「数据就绪」的信号；消费者读取之前调用 \co@wait full@，阻塞到信号到达。对称地，消费者读完（缓冲区可以复用了）就触发 \co@empty@，生产者在写下一个分块之前等待 \co@empty@。

这是一种\textbf{通过许可交接的有界缓冲区}协议——\co@full@ 是「数据可用」许可，\co@empty@ 是「缓冲区空闲」许可。


\subsection*{同步实践：流水线、缓冲区与事件}

缓冲流水线有两项顺序要求：消费者必须等生产者填完槽位后才能读取，生产者也必须等消费者用完槽位后才能覆盖。装入再快也不能省去第二项要求，消费者再快也不能省去第一项要求。

第 3.6.2 节的双缓冲模式由一个组发起搬运、等待当前输入并轮换句柄。生产者—消费者流水线则将这些职责分配给不同异步区域，再用事件传递就绪和复用信息；在单组版本中，这些顺序由组内指令安排表达。

两种模式解决的是同一个根本问题：\textbf{不能在一方写入缓冲区的同时读取它}。区别只在于解法的组织方式。


\subsection*{为什么没有事件就无法「直接重叠」}

\co@swap@ 模式之所以能够工作，是因为同一个线程组同时控制装载和计算——它清楚顺序。一旦拆成两组，就没法共享同一套 \co@swap@ 调度了，需要额外的通信机制。

下图对比了严格的「先装载再计算」阶梯和双缓冲重叠：逻辑上的工作量一样，空闲时间却少了。

图~\ref{fig:f08-fig1_pipeline_timeline} 把两种调度画在同一条时间轴上。

\begin{figure}[htbp]\centering
\includegraphics[width=\linewidth,height=0.56\textheight,keepaspectratio]{f08-fig1_pipeline_timeline_light.png}
\caption{同样的四个 K 分块可以顺序执行，也可以通过双缓冲重叠执行。图中时间仅为示意，并包含流水线启动与收尾。}
\label{fig:f08-fig1_pipeline_timeline}
\end{figure}


\subsection*{让它异步}

和 \co@parallel-by@ 语句类似，也可以让执行发散代码的线程与不执行发散代码的线程彼此异步。下面的例子展示如何在 SPMD 代码里做 MPMD 执行：

\begin{lstlisting}[language=croqtile,caption={让它异步}]
__co__  void foo() {
  parallel p by 6 {
    inthreads.async (p <= 2) {
      // divergent code for thread 0, 1, 2 (thread-mask: 000111)
    }
    inthreads.async (p > 2) {
      // divergent code for thread 3, 4, 5 (thread-mask: 111000)
    }
    sync.barrier : thread; // sync point
                 // all 'inthreads' paths must get executed
  }
}
\end{lstlisting}

本例把并行线程平均分成两个子组，并且不在每个 \co@inthreads@ 块之后强制同步。这样两组的线程就能并行执行，模拟出 MPMD 的效果。不过 \co@sync.barrier@ 语句为所有线程建立了一个\textbf{同步点}，保证两条路径的代码都执行到该\emph{同步点}之后。

只有外层 \co@inthreads@ 能声明为异步。对任何内层 \co@inthreads@ 块使用 \co@inthreads.async@，都会得到编译时错误。


\subsection*{事件实例}

在 Croqtile 中，事件与存储说明符关联，可以是 \co@global@、\co@shared@ 或 \co@local@；这个存储说明符决定了事件的实例。

\begin{lstlisting}[language=croqtile,caption={事件实例}]
__co__ void foo() {
  global event ge;     // single instance in the tileflow program
  parallel p by 2 {    // shared level
    shared event se;   // two instances, 1 for each shared instance
    parallel q by 6 {  // local level
      local event le;  // 12 instances in total, 6 for each shared instance
    }
  }
}
\end{lstlisting}

上面代码里的注释说明了事件实例。机制和带形状数据的声明相同，后者同样要求存储说明符。从实现角度看，事件可以做成指定存储层次里的布尔数据，因此也占用该层次的存储；但在概念上，事件属于某个特定层次，实例数就对应它的并行层次。


\subsection*{死锁}

默认情况下，事件是无锁的，用于不同异步块之间的高性能同步。这意味着多个线程同时更新事件时有风险，尤其是 \co@wait@。

例如：

\begin{lstlisting}[language=croqtile,caption={单实例死锁}]
parallel p by 1 {
  shared event e;
  trigger e;   // one thread set e;
  parallel q by 2 {
    wait e;    // two threads can reset e;
  }
}
\end{lstlisting}

这段代码里，两个线程等待同一个事件 \co@e@。但 \co@wait@ 会复位事件 \co@e@，因此可能出现 \co@thread-0@ 在 \co@thread-1@ 发现 \co@e@ 已触发之前就把它复位的情况，结果 \co@thread-1@ 永远无法解除阻塞，形成死锁。

\begin{lstlisting}[language=croqtile,caption={信用计数死锁}]
qc = 2;
parallel p by 1 {
  shared event e[qc];
  trigger e;   // one thread set e[0] and e[1];
  parallel q by qc {
    wait e[q]; // no dead lock
  }
}
\end{lstlisting}

这段代码使用事件数组，保证每个事件只由一个线程等待和复位。



\subsection{Canonical Pattern: One Producer and One Consumer}
\label{sec:3-8-1}

\subsection*{典型模式：1 生产者 + 1 消费者}

\co@inthreads.async@ 最常见的用途，是矩阵乘法里的 \textbf{1P1C}（一生产者、一消费者）划分：

\begin{lstlisting}[language=croqtile,caption={典型模式：1 生产者 1 消费者}]
parallel p1 by 2 : thread {
  inthreads.async (p1 == 0) {
  }
  inthreads.async (p1 == 1) {
  }
}
\end{lstlisting}

声明 \co@parallel p1 by 2 : thread@ 在同一个外层块中建立两个线程角色。\co@p1@ 的值选择职责，并不表示 K 方向的分块编号。

\co@p1 == 0@ 的区域承担生产者职责，负责写入共享暂存缓冲区。若写入使用另一个已空闲槽位，另一组就可以同时推进计算。

\co@p1 == 1@ 的区域承担消费者职责，读取已完成的分块并更新累加器。两个区域共享暂存存储，因此事件协议必须明确槽位何时从可写变为可读，以及何时重新可写。

共享地址空间要求明确的读写协调：没有协调的话，消费者可能在生产者还没写完缓冲区时就开始读。事件用于表达这一先后关系。



\subsection{Worked Example: A 1P1C Kernel Skeleton}
\label{sec:3-8-2}

\subsection*{1P1C 矩阵乘法骨架}

这个骨架省略了基于事件的同步；3.8.5 节会补上完整的流水线协议。

\begin{lstlisting}[language=croqtile,caption={一个 1P1C 矩阵乘法骨架}]
#define WARP_M 64
#define WARP_N 64
#define TILE_K 128
#define WARP_K 16
__co__ auto matmul(global f16 [M, K] lhs, global f16 [N, K] rhs) {
  f16 [M, N] output;
  parallel {block_m, block_n} by [cdiv(M, WARP_M), cdiv(N, WARP_N)] : block {
    shared f16 [WARP_M, TILE_K] lhs_load_s;
    shared f16 [WARP_N, TILE_K] rhs_load_s;
    shared f16 [WARP_M, WARP_N] output_s;
    parallel p1 by 2 : thread {
      inthreads.async (p1 == 0) {
        foreach {iv_k} in [cdiv(K, TILE_K)] {
          dma.copy lhs.subspan(WARP_M, TILE_K).at(block_m, iv_k) => lhs_load_s;
          dma.copy rhs.chunkat(block_n, iv_k) => rhs_load_s;
        }
      }
      inthreads.async (p1 == 1) {
        mc = mma.fill.f16 0.0f;
        foreach {iv_k} in [cdiv(K, TILE_K)] {
          foreach {iv_warp} in [cdiv(TILE_K, WARP_K)] {
            ma = mma.load lhs_load_s.chunkat(_, iv_warp);
            mb = mma.load rhs_load_s.chunkat(_, iv_warp);
            mma.row.row mc, ma, mb;
          }
        }
        mma.store mc, output_s;
        dma.copy output_s => output.subspan(WARP_M, WARP_N).at(block_m, block_n);
      }
    }
  }
  return output;
}
\end{lstlisting}

生产者循环遍历 \co@cdiv(K, TILE_K)@ 个输入切片，每次填充左右暂存缓冲区。对应收缩开始之前，消费者需要确认两个操作数都已就绪。

消费者为输出分块初始化一个累加器，然后针对各 K 切片装入并相乘操作数片段，在最后一次贡献加入后写回输出。两个循环独立运行，仅仅迭代次数相同，并不能保证消费者读到预期切片。

这段骨架省略了两个循环之间的交接。作为实际流水线使用之前，需要在装入后增加就绪信号，在消费后增加复用信号。第 3.8.5 节通过每槽位事件补齐这两项要求。



\subsection{Persistent Scheduling}
\label{sec:3-8-3}

前面的核函数里，网格随问题规模增长：大致每个输出分块对应一个块。矩阵一大，启动的块数就很多，而最后一批块常常让一部分 SM 闲着——这就是\textbf{尾部利用率不足}。

\textbf{持久核函数}把启动规模固定下来（通常接近 SM 数量），让每个块迭代处理多个分块：

\begin{lstlisting}[language=croqtile,caption={持久调度与 if 守卫}]
#define WARP_M 64
#define WARP_N 64
#define TILE_K 128
#define WARP_K 16
#define NUM_GROUPS 4
__co__ auto matmul(global f16 [M, K] lhs, global f16 [N, K] rhs) {
  f16 [M, N] output;
  int total_tiles = cdiv(M, WARP_M) * cdiv(N, WARP_N);
  parallel block_id by NUM_GROUPS : block {
    shared f16 [WARP_M, TILE_K] lhs_load_s;
    shared f16 [WARP_N, TILE_K] rhs_load_s;
    shared f16 [WARP_M, WARP_N] output_s;
    foreach {tile_iter} in [cdiv(total_tiles, NUM_GROUPS)] {
      tile_id = tile_iter # block_id;
      if (tile_id < total_tiles) {
        block_m = tile_id / cdiv(N, WARP_N);
        block_n = tile_id % cdiv(N, WARP_N);
        mc = mma.fill.f16 0.0f;
        foreach {iv_k} in [cdiv(K, TILE_K)] {
          dma.copy lhs.subspan(WARP_M, TILE_K).at(block_m, iv_k) => lhs_load_s;
          dma.copy rhs.subspan(WARP_N, TILE_K).at(block_n, iv_k) => rhs_load_s;
          foreach {iv_warp} in [cdiv(TILE_K, WARP_K)] {
            parallel p by 1 : thread {
              ma = mma.load lhs_load_s.chunkat(_, iv_warp);
              mb = mma.load rhs_load_s.chunkat(_, iv_warp);
              mma.row.row mc, ma, mb;
            }
          }
        }
        mma.store mc, output_s;
        dma.copy output_s => output.subspan(WARP_M, WARP_N).at(block_m, block_n);
      }
    }
  }
  return output;
}
\end{lstlisting}

声明将工作块数量固定为 \co@NUM_GROUPS@。工作量变化时，调整的是循环迭代次数，而不是启动块数。

表达式 \co@tile_iter # block_id@ 的值为 \co@tile_iter * NUM_GROUPS + block_id@，每次迭代为各工作块分配下一组连续分块。

条件 \co@tile_id < total_tiles@ 排除最后一个有效分块之后的索引。属于该分块的工作都应放入判断之内，包括需要参与者一致确认分块是否存在的生产者—消费者协议。

图~\ref{fig:f07-fig2_persistent_kernel} 展示了十个分块如何分给四个固定块。

\begin{figure}[htbp]\centering
\includegraphics[width=\linewidth,height=0.56\textheight,keepaspectratio]{f07-fig2_persistent_kernel_light.png}
\caption{四个固定 block 用三轮迭代处理十个分块；最后两个补齐的分块编号由 if 条件跳过。}
\label{fig:f07-fig2_persistent_kernel}
\end{figure}

若有四个工作块和十个输出分块，组合索引在第一次迭代为 0–3，第二次为 4–7，第三次为 8–11。最后一次只有 8 和 9 通过条件判断。这样每个有效分块恰好被处理一次，但并不保证所有工作块始终忙碌，利用率仍取决于分块开销和剩余工作量。


\subsection*{数据相关网格与持久网格}

\begin{table}[htbp]\centering\small
\begin{tabularx}{\linewidth}{@{}lXX@{}}
\toprule
\textbf{方面} & \textbf{每分块一个块} & \textbf{持久（\texttt{NUM\_GROUPS} 块）} \\
\midrule
网格规模 & 随问题增长 & 固定 \\
尾部利用率 & 最后一批可能使 SM 空闲 & 可继续处理后续分块；尾部仍可能空闲 \\
额外构造 & 最少 & \texttt{total\_tiles}、\texttt{tile\_iter \# block\_id}、\texttt{if} \\
复杂度 & 较低 & 较高 \\
\bottomrule
\end{tabularx}
\end{table}



\subsection{Host-Level Concurrency and Stream Binding}
\label{sec:3-8-4}

\subsection*{\texttt{parallel.async}：主机层并发}

上面这些都发生在核函数内部。有时需要在\textbf{主机层}实现并发：启动一个网格而不阻塞 CPU。

\begin{lstlisting}[language=croqtile,caption={parallel.async：主机层并发}]
parallel.async {px, py} by [grid_m, grid_n] : block {
  // kernel body
}
\end{lstlisting}

异步启动在提交核函数后将控制权返回主机，完成发生在之后。流绑定选择接收工作的流，没有显式绑定时使用默认流。

这属于\textbf{主机编排}，与核函数内部的控制流互不相干。它既不能替代用于线程划分的 \co@inthreads.async@，也不能替代用于运行时谓词的 \co@if@——它决定的是一个网格相对于其他网格\emph{在何时、何处}运行。


\subsection*{异步程序启动与流绑定}

默认情况下，主机会等 \co@parallel-by@ 核函数完成后才继续。要异步启动核函数，就用 \co@parallel.async@。

\begin{lstlisting}[language=croqtile,caption={异步程序启动与流绑定}]
__co__ void foo(stream s, f32[M] a) {
  parallel.async<s> p by M : block {
    a.at(p) = 1.0f;
  }
  // host continues immediately (no synchronization)
}
\end{lstlisting}

3.8.4 节对此有详细说明。


\subsection*{同步与异步 \texttt{parallel-by}}

和 \co@inthreads@ 类似，\co@parallel-by@ 也可以加上 \co@.async@ 变成异步：

\begin{lstlisting}[language=croqtile,caption={一次异步启动}]
parallel.async p by 6 : block {
}
\end{lstlisting}

不加 \co@.async@ 时，主机会等核函数结束后才继续：

\begin{lstlisting}[language=croqtile,caption={同步形式}]
parallel p by 6 : block {
}
// host continues only after the kernel finishes
\end{lstlisting}

加上 \co@.async@ 之后，主机代码在启动核函数后立刻继续。要让计算和数据传输重叠，或者并发启动彼此独立的核函数，这一点很关键。


在 Croqtile 中，可以用尖括号语法把 \co@parallel-by@ 块绑定到指定的流：

\begin{lstlisting}[language=croqtile,caption={流绑定}]
__co__ void foo(stream s0, stream s1, f32[M] a, f32[N] b) {
  parallel.async<s0> p by M : block {
    a.at(p) = 1.0f;
  }
  parallel.async<s1> q by N : block {
    b.at(q) = 2.0f;
  }
}
\end{lstlisting}

这个例子里：

\begin{itemize}[leftmargin=1.6em]
\item **\co@stream s0, stream s1@**：由主机传入的流参数。
\item **\co@parallel.async<s0>@**：在流 \co@s0@ 上异步启动核函数。
\item **\co@parallel.async<s1>@**：在流 \co@s1@ 上启动第二个核函数。
\end{itemize}


\subsection*{不带 \texttt{.async} 的流绑定}

流绑定同样适用于同步的 \co@parallel-by@。此时主机在指定的流上启动核函数，然后等待该流完成：

\begin{lstlisting}[language=croqtile,caption={不带 .async 的流绑定}]
__co__ void bar(stream _s, f32[M] a) {
  parallel<_s> p by M : block {
    a.at(p) = 1.0f;
  }
  // host synchronizes on stream _s before continuing
}
\end{lstlisting}


\subsection*{规则}

\begin{itemize}[leftmargin=1.6em]
\item 流绑定只能用在\textbf{块级}并行（\co@: block@）上。
\item \co@<...>@ 里的标识符必须是 \co@stream@ 类型的变量。
\item 不做流绑定时，同步的 \co@parallel-by@ 使用设备范围的同步机制。
\end{itemize}

异步的 \co@parallel-by@ 则使用默认流。

流为启动操作选择顺序上下文。如果第二个核函数读取第一个写入的数据，使用流不会让两者变得独立。使用不同流时，仍需明确表达依赖，并将缓冲区保留到队列中的使用者结束。异步启动返回主机，只说明执行已提交。



\subsection{Worked Example: A Full Event Pipeline}
\label{sec:3-8-5}

\subsection*{完整的 1P1C 事件流水线}

这个核函数把 3.8 节的 \co@inthreads.async@ 和事件数组结合起来，构建一条完整的多级流水线。生产者和消费者作为独立的并发程序运行，通过 \co@wait@ / \co@trigger@ 协调：

\begin{lstlisting}[language=croqtile,caption={完整的 1P1C 事件流水线}]
#define TILE_M 128
#define TILE_N 128
#define TILE_K 128
#define WARP_M 64
#define WARP_N 64
#define WARP_K 16
#define STAGES 2
__co__ auto matmul(global f16 [M, K] lhs, global f16 [N, K] rhs) {
  f16 [M, N] output;
  parallel {block_m, block_n} by [cdiv(M, WARP_M), cdiv(N, WARP_N)] : block {
    shared event full[STAGES], empty[STAGES];
    shared f16 [WARP_M, TILE_K] lhs_load_s[STAGES];
    shared f16 [WARP_N, TILE_K] rhs_load_s[STAGES];
    shared f16 [WARP_M, WARP_N] output_s;
    parallel p1 by 2 : thread {
      inthreads.async (p1 == 0) {
        foreach {iv_k} in [cdiv(K, TILE_K)] {
          stage = iv_k % STAGES;
          wait empty[stage];
          dma.copy lhs.subspan(WARP_M, TILE_K).at(block_m, iv_k)
            => lhs_load_s[stage];
          dma.copy rhs.chunkat(block_n, iv_k)
            => rhs_load_s[stage];
          trigger full[stage];
        }
      }
      inthreads.async (p1 == 1) {
        mc = mma.fill.f16 0.0f;
        foreach {s} in [STAGES]
          trigger empty[s];
        foreach {iv_k} in [cdiv(K, TILE_K)] {
          stage = iv_k % STAGES;
          wait full[stage];
          foreach {iv_warp} in [cdiv(TILE_K, WARP_K)] {
            ma = mma.load lhs_load_s[stage].chunkat(_, iv_warp);
            mb = mma.load rhs_load_s[stage].chunkat(_, iv_warp);
            mma.row.row mc, ma, mb;
          }
          trigger empty[stage];
        }
        mma.store mc, output_s;
        dma.copy output_s => output.subspan(WARP_M, WARP_N).at(block_m, block_n);
      }
    }
  }
  return output;
}
\end{lstlisting}


\subsection*{逐段走查核函数}

阶段索引为 \co@iv_k % STAGES@。使用两个槽位时，分块 0 和 2 使用槽位 0，分块 1 和 3 使用槽位 1。它们是不同的逻辑分块，在不同时刻共用物理存储；事件对属于槽位，不属于某个永久分块。

生产者在写入任一操作数前等待 \co@empty[stage]@。本例的两次复制均为同步操作，因此在它们之后触发 \co@full[stage]@，就表示两个操作数都已就绪。若替换为异步搬运，则必须先完成这些搬运，再发出同一个触发。

消费者先触发所有 \co@empty@ 事件。事件初始均未置位，因此这一步授予生产者首次填充各槽位的许可。此时尚未装入输入，所以不会触发任何 \co@full@ 事件。

消费者等待 \co@full[stage]@，装入操作数片段并累加当前切片，随后触发 \co@empty[stage]@。这个顺序不能颠倒：发出矩阵操作，并不足以证明它已经停止读取共享操作数。

完成点保护缓冲区复用。若 MMA 仍在读取时生产者就覆盖槽位，即使分块形状完全正确，消费者也可能混用两个 K 切片的操作数。仅检查 \co@full@ 事件无法防止这一错误，归还 \co@empty@ 事件同样必要。


\subsection*{单级的事件流}

图和代码一一对应：引导阶段先触发每个 \co@empty@ 事件；生产者等待 \co@empty@、填入数据、触发 \co@full@；消费者等待 \co@full@、做计算、触发 \co@empty@。当 \co@iv_k@ 按 \co@STAGES@ 取模回绕时，同一个物理级又进入下一轮循环。

图~\ref{fig:f08-fig2_event_credit_flow} 跟踪了两个事件在单级流水中的流转。

\begin{figure}[htbp]\centering
\includegraphics[width=\linewidth,height=0.56\textheight,keepaspectratio]{f08-fig2_event_credit_flow_light.png}
\caption{\co@empty[s]@ 允许生产者填充该槽位，\co@full[s]@ 允许消费者读取；消费者在计算完成后才释放该阶段供再次使用。}
\label{fig:f08-fig2_event_credit_flow}
\end{figure}



\subsection{Coordination Pattern Summary}
\label{sec:3-8-10}

\subsection*{小结}

\begin{table}[htbp]\centering\small
\begin{tabularx}{\linewidth}{@{}lXX@{}}
\toprule
\textbf{概念} & \textbf{原语} & \textbf{何时使用} \\
\midrule
谓词执行 & \co@if@ & 数据相关的判断（边界、条件） \\
结构化并发 & \co@inthreads.async@ & 编译时划分线程（生产者／消费者、异构角色） \\
区域间通信 & \co@shared event@ / \co@wait@ / \co@trigger@ & 并发区域之间的协调 \\
主机并发 & \co@parallel.async@ / \co@stream s@ & 多核函数重叠、非阻塞启动 \\
持久调度 & \co@if@ + \co@foreach@ + \co@#@ & 固定网格规模、带填充守卫的分块条带化 \\
\bottomrule
\end{tabularx}
\end{table}

完整流水线为前面的骨架补上两项依赖：计算前输入必须就绪，复用槽位前操作数必须用完。修改流水线时，可以沿一个槽位追踪完整周期，再追踪下一次复用。与只看某一侧循环相比，这样更容易发现缺失信号和迭代次数不匹配。


\subsection*{小结}

\begin{table}[htbp]\centering\small
\begin{tabularx}{\linewidth}{@{}lXX@{}}
\toprule
| \textbf{模式} & \textbf{所用原语} & \textbf{何时使用} | \\
\midrule
| --- & --- & --- | \\
| \co@swap@ 双缓冲 & \co@dma.any@、\co@swap@ & 单线程组交错执行装载与计算 | \\
| 事件流水线 & \co@inthreads.async@、事件 & 生产者与消费者分别推进 | \\
\bottomrule
\end{tabularx}
\end{table}

一个组同时控制搬运调度和计算时，可以通过句柄轮换组织缓冲。分离的生产者与消费者需要事件，因为任一侧都无法仅凭自身指令顺序判断另一侧进度。增加缓冲区可以为重叠执行提供空间，却无法修复缺失的依赖。

更深的流水线则改用事件数组与 \co@inthreads.async@，底层数据搬运原语不变。



\section{C++ 互操作与预处理器}
\label{sec:3-9}

\subsection*{C++ 互操作：设备函数、内联代码与预处理器}

核函数可能需要归约、排序过程，或 Croqtile 没有直接表达的指令。已有设备 C++ 可以提供这部分计算，Croqtile 继续安排输入、输出和并行范围。调用接口必须保留两侧代码各自的使用约定。

Croqtile 的互操作方案由三部分组成：

\begin{itemize}[leftmargin=1.6em]
\item \textbf{\texttt{\_\_device\_\_} 函数} —— 编译器保留原样。
\item \textbf{\texttt{\_\_cpp\_\_}}
\item \textbf{预处理器} —— \co@#define@、\co@#if@、\co@#ifdef@、\co@#error@ 用于编译时配置，作用与 C/C++ 代码库中的 \co@#define@ 相同。
\end{itemize}

前面几节介绍了分块、并行、MMA、流水线、事件等抽象。需要额外的设备函数或目标平台代码时，可以通过本节的 C++ 互操作机制补充。


\subsection*{\texttt{\_\_device\_\_} 与普通 C++ 内联代码}

\co@__device__@ 函数与普通 C++ 函数都会被保留原样，但二者进入不同的编译路径——这个区别很重要。

\co@__device__@ 函数编译成设备代码；\co@__co__@ 函数体用 \co@call@ 关键字调用它，生成的 \co@__global__@ 核函数会在设备上直接调用。

普通 C++ 函数（没有 \co@__device__@ 或 \co@__co__@ 标注）编译成 \textbf{CPU 代码}，运行在主机上，负责建立缓冲区、启动核函数、打印结果。不能在 \co@__co__@ 函数体里调用主机函数。

边界很严格：\co@__device__@ 在 \co@c++@ 划分的设备侧，主机 C++ 在主机侧。\co@__co__@ 函数是桥——编译器把它变换成 \co@__global__@ 核函数（设备侧），同时生成启动它的主机侧包装。


\subsection*{与 C++ 预处理的区别}

Croqtile 的预处理比 C++ 预处理早得多。Croqtile 编译流程如下图所示：

\begin{lstlisting}[language=croqtile,caption={与 C 预处理的区别}]
chore-preprocessing -> croqtile compilation -> c/c++ preprocessing -> c/c++ compilation
\end{lstlisting}

Croqtile 预处理的主要目标，是让主机和设备的宏作为一个整体生效，但这样的工作流有时会带来差异。从实现角度看，Croqtile 预处理器只对 tileflow 函数内部的代码做替换或条件编译，其余预处理留给 C++ 预处理器——这就把 Croqtile 预处理限制在了较小的范围内。


\subsection*{预定义宏}

为了模仿目标平台的原生编译，Croqtile 预处理也会取用目标平台的内建宏。因此这些宏既能在 tileflow 函数里用，也能在主机代码里用。





\subsection*{触发 C++ 模板实例化}

Croqtile 代码可以调用模板函数。例如：

\begin{lstlisting}[language=croqtile,caption={一个模板化的设备辅助函数}]
template<int M, int N, int K>
__device__ void matmul_kernel(int *lhs, int* rhs, int* output) {}
__co__ auto matmul(s32 [96, 72] lhs, s32 [72, 24] rhs) {
  s32 [lhs.span(0), rhs.span(1)] output;
  parallel {p, q} by [6, 6] : block {
    shared s32 [output.span / {#p, #q}] buffer;   // shape: [16, 4]
    lhs_load = dma.copy lhs.chunkat(p, _) => shared;
    rhs_load = dma.copy rhs.chunkat(_, q) => shared;
    call matmul_kernel<buffer.span(0), buffer.span(1), 72>(lhs_load.data, rhs_load.data, buffer);
    dma.copy buffer => output.chunkat(p, q);
  }
  return output;
}
\end{lstlisting}

调用向 \co@matmul_kernel@ 提供三个常量模板实参，目标编译器据此实例化对应特化，与普通 C++ 模板调用相同。

所以，在 Croqtile 中调用模板函数与 C++ 类似。但\textbf{传入的模板实参必须能被 Croqtile 编译器推导为编译时常量}；带运行时值时，编译器会给出警告，指出该实参无法在编译时求值。

\begin{lstlisting}[language=croqtile,caption={一次触发实例化的调用}]
__co__ void foo(int M) {
  parallel p by 1 {
    call bar<M>();  // warning: 'M' cannot be evaluated at compile time
  }
}
\end{lstlisting}

Croqtile 编译器会尽可能在编译时推导取值。若模板实参无法推导，编译器会给出警告，指出该实参无法在编译时求值。


\subsection*{更多问题}

和 C++ 函数一样，Croqtile 函数可以定义在源代码的任何位置。把模块和 Croqtile 函数定义链接起来时，需要确定它的\textbf{函数签名}。

答案是：Croqtile 函数遵循 C++ 调用约定，也就是带 C++ 名字改编的标准 C++ 函数。Croqtile 函数\textbf{不能}使用 \co@extern "C"@，因为它的参数类型（带形状数据）是基于 C++ 模板的对象。

还有一些特性尚未支持：

\begin{itemize}[leftmargin=1.6em]
\item 只有声明、没有定义的 Croqtile 函数——有了它，才能从另一个模块调用 Croqtile 函数。
\item 带 \co@inline@ 的 Croqtile 函数——有了它，Croqtile 函数才能放进头文件。
\end{itemize}



\subsection{Worked Example: Calling a Device Helper}
\label{sec:3-9-1}

\subsection*{\texttt{\_\_device\_\_} 函数：与 Croqtile 一并面向目标平台}

自定义归约、排序网络和特殊数学运算等计算，可以通过设备辅助函数实现。

\co@__co__@ 函数可以用 \co@call@ 关键字调用 \co@__device__@ 函数：

\begin{lstlisting}[language=croqtile,caption={device 函数：与 Croqtile 一并面向目标平台}]
#define N_TOKENS 1024
#define N_EXPERTS 8
#define N_THREAD 128
template <int K>
__device__ void warp_topk(f32* vals, s32* idxs);
template <typename T>
__device__ __forceinline__ T SHFL_XOR(T var, int lane_mask, int width) {
  return __shfl_xor_sync(uint32_t(-1), var, lane_mask, width);
}
__co__ void moe_topk(f32 [N_TOKENS, N_EXPERTS] scores,
                     s32 [N_TOKENS, K]& topk_ids,
                     f32 [N_TOKENS, K]& topk_scores,
                     int N_BLOCK) {
  parallel n by N_BLOCK : block {
    foreach m in [ |scores.span| / N_THREAD / N_BLOCK ] {
      shared_scores = dma.copy scores.chunkat(n#m, _) => shared;
      parallel gid by [N_THREAD / 32] : group
        parallel tid by 32 : thread {
          flat_scores = shared_scores.span_as([|shared_scores.span|]);
          score = flat_scores.at(gid # tid);
          local s32 [K] frag_idx{-1};
          local f32 [K] frag_val{-1.0f};
          frag_idx.at(0) = tid;
          frag_val.at(0) = score;
          call warp_topk<8>(frag_val, frag_idx);
          inthreads (tid == 0)
            foreach k in K {
              topk_ids.at(n#m, gid#k) = frag_idx.at(k);
              topk_scores.at(n#m, gid#k) = frag_val.at(k);
            }
        }
    }
  }
}
\end{lstlisting}

辅助函数接收指向元素存储的原始指针。因此，接口需要给出它所需的尺寸或布局约定，仅凭指针不会保留 Croqtile 形状。

\co@SHFL_XOR@ 辅助函数是交给目标编译器的普通设备代码。计算中如何使用它属于辅助函数的实现，Croqtile 不会把其函数体翻译成分块操作。

语句 \co@call group_topk<8>(...)@ 提供辅助函数的编译时配置和普通实参。Croqtile 将带形状数据实参转换为元素指针，由目标编译器验证生成的 C++ 调用。

区域 \co@inthreads (tid == 0)@ 选择一个线程完成其中工作，结束后汇合。它没有 \co@.async@ 修饰符，因此不会引入独立推进的生产者或消费者。

这种分工——Croqtile 管并行、分块和内存编排，设备函数管专门计算——在 MoE（专家混合）top-K、自定义归约和专用数学运算等生产核函数里很常见。



\subsection{Interoperability Summary}
\label{sec:3-9-10}

\subsection*{小结}

\begin{table}[htbp]\centering\small
\begin{tabularx}{\linewidth}{@{}lXX@{}}
\toprule
| \textbf{语法} & \textbf{作用} & \textbf{运行于} | \\
\midrule
| --- & --- & --- | \\
| \co@#define NAME value@ & 对象式宏；在 \co@__co__@ 与主机区域间共享 & 预处理器 | \\
| \co@#if@ / \co@#ifdef@ / \co@#error@ & 条件编译与编译时断言 & 预处理器 | \\
| \co@-DNAME=value@ & 从命令行定义宏，用于分块尺寸扫描 & 预处理器 | \\
| \co@call helper(...)@ & 调用设备辅助函数 & 设备 | \\
| \co@__cpp__("...")@ & 在当前位置插入目标 C++ & 取决于所在目标作用域 | \\
\bottomrule
\end{tabularx}
\end{table}

3.11 节转向工作流程的另一面：核函数能编译、输出却不对时该怎么办——调试、详细模式，以及一步步缩小范围。3.10 节的两个完整程序，则把互操作机制放进了具体语境。



\section{Tileflow 模式与设计依据}
\label{sec:3-10}

\subsection*{受限接口为何有效：表达力与性能}

逐条编写装入和地址计算给程序员很多选择，也让编译器需要重新识别哪些操作共同构成一次分块搬运。Croqtile 要求程序显式描述这片区域。当所需搬运能够用规则形状表示时，这项约束很有用，因为在生成各次传输之前，区域几何信息就已明确。

Croqtile 用结构化搬运操作显式描述分块形状和布局。编译器据此生成传输代码，程序员也不必为每个元素反复编写地址计算与协作逻辑。

图~\ref{fig:f09-fig8_expressiveness} 对比了「能写出来的范围」与「跑得快的范围」。

\begin{figure}[htbp]\centering
\includegraphics[width=\linewidth,height=0.56\textheight,keepaspectratio]{f09-fig8_expressiveness_light.png}
\caption{结构化分块操作显式提供形状和布局，便于编译器生成高效代码；实际性能仍取决于分块选择与执行配置。}
\label{fig:f09-fig8_expressiveness}
\end{figure}

这里的取舍是用明确的分块表达替代任意的逐元素搬运代码，从而便于高效生成目标代码。实际性能仍取决于分块尺寸、对齐方式、布局和同步安排，并不由语法单独保证。


\subsection*{易于使用}

分块选择可以直接使用输入形状和有界索引。例如，改变分块数量后，派生分块尺寸随之改变，无需另行维护一套手工偏移计算。划分仍由程序员选择，语言则把这一选择带来的相关信息连接起来。


\subsection*{编译时安全}

信息充分时，形状和秩检查能在执行前发现不兼容的用法。运行时检查覆盖依赖实际输入的约束，断言则可补充应用自身的要求。这些检查有助于定位失败，但不能替代输出数值验证和同步协议检查。


\subsection*{动态形状}

符号维度让一份源码表达一组输入形状。输出大小和迭代次数可以从命名维度推导，不必作为独立常量重复维护。每次调用仍需满足核函数的分块、布局和目标约束，符号形状并不意味着所有形状都能得到同样好的支持。


\subsection*{为 AI 调优而生}

紧凑源码让自动调优过程能够同时保留核函数、编译诊断和测量结果。显式分块尺寸与协作范围也提供了可调整的具体选择。每个候选仍需验证正确性，并在目标负载上测量；成功编译只是其中一部分依据。


\subsection*{Croqtile 中的分块优化}

下面的示例将相同分块选择机制用于不同计算。矩阵乘法需要在 K 切片之间保留部分和，RGB 转灰度则独立处理每片图像区域。对照数据路径，可以看出哪些结构来自分块安排，哪些来自计算本身。



\subsection{Worked Example: Tiled matmul}
\label{sec:3-10-1}

\subsection*{示例 1：\texttt{matmul}}

第一个例子是设备程序上的矩阵乘法。代码展示了 Croqtile 的分块和内存管理机制如何配合起来优化矩阵运算，也包括如何借助 DMA（直接内存访问）高效搬运数据。


\subsection*{示例代码：带分块的矩阵乘法}

\begin{lstlisting}[language=croqtile,caption={示例代码：带分块的矩阵乘法}]
__co__ s32 [512, 1024] matmul(s32 [512, 1024] lhs, s32 [1024, 1024] rhs) {
  // Device program
  s32[lhs.span(0), rhs.span(1)] output; // Use same shape as lhs
  parallel p by 1 {
    with index = {m_tile, k_tile, n_tile} in [4, 4, 8] {
      foreach m_tile, n_tile {
        shared s32[128, 128] l2_out;
        foreach k_tile {
          // Initiate DMA operations
          lhs_load = dma.copy.async lhs.chunkat(m_tile, k_tile) => shared;
          rhs_load = dma.copy.async rhs.chunkat(k_tile, n_tile) => shared;
          out_store_s = dma.any;
          parallel q by 1 {
            with index = {m_tile_s, k_tile_s, n_tile_s} in [4, 4, 8] {
              // Local buffers for tiling
              local s32 [32, 64] lhs_load_s_buffer_0 {0}, lhs_load_s_buffer_1 {0};
              local s32 [64, 16] rhs_load_s_buffer_0 {0}, rhs_load_s_buffer_1 {0};
              foreach index {
                // Add "after" primitive to chain up paired DMA operations
                lhs_load_s = dma.copy.async lhs_load.data.chunkat(m_tile_s, k_tile_s) => select(k_tile_s % 2, lhs_load_s_buffer_0, lhs_load_s_buffer_1) after lhs_load;
                rhs_load_s = dma.copy.async rhs_load.data.chunkat(k_tile_s, n_tile_s) => select(k_tile_s % 2, rhs_load_s_buffer_0, rhs_load_s_buffer_1) after rhs_load;
                out_load_s = dma.copy.async l2_out.chunkat(m_tile_s, n_tile_s) => local;
                wait out_load_s;
                // Perform the matrix multiplication (dot product) operation
                wait lhs_load_s, rhs_load_s;
                call dot_general_fp32(lhs_load_s.data, rhs_load_s.data, 32, 16, 64, 0, 8, out_load_s.data);
                // Store the result back to l2_out
                out_store_s = dma.copy.async out_load_s.data => l2_out.chunkat(m_tile_s, n_tile_s);
              }
            }
          }
          // Add "after" primitive to chain up DMA operations for the final output
          out_store = dma.copy.async l2_out => output.chunkat(m_tile, n_tile) after out_store_s;
          wait out_store;
        }
      }
    }
  }
  return output;
}
\end{lstlisting}

外层范围 \co@[4, 4, 8]@ 把操作数切成 \co@[128, 256]@ 的左分块、\co@[256, 128]@ 的右分块，以及 \co@[128, 128]@ 的输出分块——这正是 \co@l2_out@ 暂存缓冲区的大小。内层同样按 \co@[4, 4, 8]@ 划分，得到 \co@[32, 64]@、\co@[64, 16]@ 和 \co@[32, 16]@ 的结果。\co@after@ 子句使内层输入复制在外层复制之后开始。这段清单用于展示结构；作为可运行核函数使用前，必须补齐并检查辅助函数定义、操作数使用与缓冲区复用前的完整等待，以及输出初始化。


\subsection*{用 \texttt{chunkat} 做内存分块}

矩阵乘法示例里，\textbf{内存分块}是用 \co@chunkat@ 方法做的：它把输入矩阵 \co@lhs@（左操作数）和 \co@rhs@（右操作数）切成更小的块，方便并行处理。

\begin{itemize}[leftmargin=1.6em]
\item \co@lhs@ 与 \co@rhs@ 矩阵用 \co@chunkat@ 方法分块，分块由 \co@m_tile@、\co@k_tile@、\co@n_tile@ 定义，分别对应乘法涉及矩阵的不同维度。
\item 两个矩阵的分块都用 \co@dma.copy.async@ 异步装入（\co@lhs_load@ 和 \co@rhs_load@）。这个操作把矩阵的一部分从源缓冲区搬进暂存缓冲区，是大型矩阵运算里改善数据局部性的常见优化。
\end{itemize}

\begin{itemize}[leftmargin=1.6em]
\item 变量 \co@m_tile@、\co@k_tile@、\co@n_tile@ 是\textbf{有界变量}，定义在 \co@parallel-by@ 和 \co@with-in@ 构造的上下文中，这些构造确定了分块的迭代空间。
\item 这些有界变量让分块能被并行计算，每个线程处理矩阵的不同切片，矩阵乘法因此能跨多个线程高效完成。
\end{itemize}



\subsection{Worked Example: rgb2gray}
\label{sec:3-10-2}

\subsection*{示例 2：rgb2gray}

第二个例子是 RGB 转灰度算法：用加权平均把彩色 RGB 图像变成灰度图像。它展示了物理并行和虚拟并行如何协同工作。


\subsection*{示例代码}

\begin{lstlisting}[language=croqtile,caption={示例代码}]
__co__ auto rgb2gray(f32 [N, 3, H, W] input) {
  f32 [N, H, W] out;
  parallel q by 6 {
    with index={n, h, w} in [N, H, W]/{#q, 16, 512} {
      foreach n, h, w {
        // _ means no tiling
        input_L1_A = dma.copy input.chunkat(q#n, _, h, w) => local;
        dims : input_L1_A.span[(0), (2), (3)];
        local f32 [dims] out_L1;
        call rgb2gray_kernel_fp32_fp32(input_L1_A.data, out_L1, |out_L1|, 1);
        dma.copy out_L1 => out.chunkat(q#n, h, w);
      }
    }
  }
  return out;
}
\end{lstlisting}

输入形状为 \co@[N, 3, H, W]@，第二维包含三个颜色通道。\co@chunkat(q#n, _, h, w)@ 中的下划线保持这些通道完整，同时选择一幅图像中的空间区域。从装入形状中去掉通道维，就得到 \co@out_L1@ 的形状，辅助函数为每个空间位置写入一个灰度值。

组合索引 \co@q # n@ 为六个并行实例分别分配一组连续图像。这种规则划分假定 N 能被 6 整除，H 能被 16 整除，W 能被 512 整除。满足这些条件时，每个输入分块包含一个 16 × 512 区域的全部三个通道，最后的复制将灰度结果写入对应输出分块。其他图像尺寸需要明确的边界处理方式。



\section{诊断与方法论}
\label{sec:3-11}

\subsection*{调试与详细模式：打印、断言与 GDB}

调试的基本方法是\textbf{逐步缩小问题范围}。先做静态分析和编译时断言，在不增加运行时开销的情况下发现错误；再用有针对性的运行时检查定位问题；范围缩小后，按需使用交互式调试器深入检查。

Croqtile 在每一层都有工具：\textbf{编译时形状打印}（\co@print!@、\co@println!@），不用启动核函数就能核对分块尺寸；\textbf{运行时断言}（\co@assert@），用来检查不变式；\textbf{带守卫的设备 \texttt{println}}，在运行时查看特定线程；给 \co@debugger@ 用的\textbf{调试 RTTI}；以及查看编译器实际执行步骤的\textbf{详细模式}（\co@-v@）。

先使用容易计算参考结果的小输入。矩阵某个元素错误时，还原它的分块坐标和块内坐标，检查选中的输入切片，再比较一个部分和，最后才看完整输出。程序挂起时，则追踪一个事件槽位，记录初始状态、每次触发以及应消费该事件的实例。这样可以区分索引错误与交接缺失。

若失败发生在编译阶段，先确定阶段再修改核函数。预处理输出显示宏展开结果，生成源码显示 Croqtile 的翻译，详细输出则给出目标编译命令。修复后，关闭诊断输出，再次与参考结果比较，然后测量性能。第 1.4 和 2.14 节给出对应命令与检视语法。



\guidechapternext{Appendix~\ref{app:quickref} collects the syntax in one place. For the compiler design behind the language, see \citep{croqtileCompiler}.}

